\quellseite{154}
\section[§ 17. The Period of Pietism]{§ 17. Religious Instruction in the Church During the Period of Pietism (About 1677–1750)}

G. Langemack, \emph{Historia catechetica}, iii, 1740. Ph. Schuler, pp. 107–207. F. Ehrenfeuchter, pp. 48–61. *G. v. Zezschwitz i, ch. 26; ii 2, 1, ch. 32; ii 2, 2, ch. 16. *Schumann-Sperber, pp. 49 ff. F. Cohrs, \emph{Katechismusunterricht} in PRE³, 1901. M. v. Nathusius, \emph{Das Ziel des kirchlichen Unterrichts}, pp. 56 ff., 1903. E. Chr. Achelis, pp. 320 ff. J. Steinbeck, pp. 43 ff. *R. Vormbaum, \emph{Evang. Schulordnungen des 18 Jahrh.}, 1864. K. v. Raumer, \emph{Geschichte der Paedagogik} ii, ⁶1897. PHIL. SPENER: *P. Gruenberg, \emph{Phil. Jac. Spener}, 3 vols., esp. ii, pp. 58–90, 1893. 1905. 1906. W. Thilo, \emph{Spener als Katechet}, 1840. W. Thilo, \emph{Der Bibelspruch im Dienst des evangelischen Religionsunterrichts}, 1846. *W. Caspari, \emph{Die evangelische Konfirmation}, esp. pp. 84 ff., 1890. W. Diehl, \emph{Zur Geschichte der Konfirmation}, pp. 90 ff., 1897. A. Ernst u. E. Adam, \emph{Katech. Gesch. des Elsasses}, p. 210, 1897. Spener’s \emph{Kurze Katechismuspredigten} (i.e. the exordia of his sermons delivered during the seventies but published later), 1689; new edition, St. Louis and Leipzig, ²1869. His *\emph{Einfältige Erklärung der christlichen Lehre nach der Ordnung des kleinen Katechismus Luthers}, 1677; new edition, Berlin, ⁵1864; translated into Tamil by Barth. Ziegenbalg, 1719, revised edition Tranquebar, 1872. \emph{Tabulae Catecheticae quibus 5 capita catechismi minoris et subnexa tabula oeconomica in certa pensa distributa . . . tractantur}, 1683; in German by J. G. Pritius, \emph{Katechismustabellen}, 1713. \emph{Gedanken von der Katechismus-Information}, aus Speners Schriften von einem Freund katechetischer Uebungen zusammengetragen, 1715. \emph{Theologische Bedenken}, 4 parts, 1700–1702. \emph{Letzte theologische Bedenken}, 1711. \emph{Consilia et judicia theologica latina}, 1709; a selection of his German and Latin \emph{Consilia} was edited by F. A. E. Hennicke, 1838. His \emph{Pia Desideria} in condensed and modernized form edited by Gruenberg in \emph{Bibliothek theolog. Klass.}, 1889. AUG. HERM. FRANCKE: G. Kramer, \emph{A. H. Francke, ein Lebensbild}, 2 vols., 1880–1882. M. E. Richard, \emph{Aug. Herm. Francke and His Work}, 1897. G. F. Herzberg, \emph{A. H. Francke und sein Hallisches Waisenhaus}, 1898. O. Schulze, \emph{Franckes Paedagogik}, 1898. *A. Otto, \emph{A. H. Francke, vol. ii: Beurteilung und Bedeutung der Franckeschen Paedagogik}, 1904. *G. Kramer, \emph{Paedagogische Schriften Franckes}, ²1885. \emph{Der}
\quellseite{155}
\emph{kurze und einfältige Bericht, wie die Kinder zur wahren Gottseligkeit und christlicher Klugheit anzuführen sind} (1702), according to the edition of 1748 re-edited by O. Frick, 1889. The two important \emph{School Regulations} are reprinted by Vormbaum, op. cit., pp. 1–116. Johann Huebner: F. Brachmann, Johann Huebner, Johannei Rector 1711–1731, 1899. G. v. Zeschwitz ii 2, 1, ch. 32. Johann Jacob Rambach: C. Bertheau, \emph{J. J. Rambach} in PRE,3 1905. *\emph{Der Wohlunterrichtete Katechet}, 1722, ⁶1730. \emph{Darmstadt School Regulations} of 1733, reprinted by Vormbaum, op. cit., pp. 343 ff. \emph{Das Erbauliche Handbuechlein fuer Kinder} (1. Order of Salvation; 2. Treasures of Salvation; 3. A little Hymn-book; 4. A new Prayer-book; 5. Stories of Pious Children; 6. Rules for a Christian Life; 7. Necessary Rules for a Good Life), 1734, ⁵1735 (!). *\emph{Der Wohlunterwiesene Informator}, containing not a few modern ideas, published after Rambach’s death by Neubauer, 1737, with a good introduction in which he shows that this work as well as \emph{Der Wohlunterrichtete Katechet} have grown forth from the catechetical and pedagogic lectures which Rambach delivered at the university of Jena. Frederick the Great and the Prussian School Regulations: J. B. Meyer, \emph{Friedrich der Grosse, paedagogische Schriften und Aeusserungen}, ²1890. Huebler, \emph{Friedrich d. Gr. als Paedagog}, 21900. E. Clausnitzer, \emph{Zur Geschichte der preussischen Volksschule unter Friedr. d. Grossen}, 1901. A. Heubaum, \emph{Geschichte des deutschen Bildungswesens} etc., esp. part 4, 1905. F. Fischer, \emph{Preussische Volksschulverordnungen}, 1907. H. Kiehl, \emph{Joh. Julius Hecker}, Jahresbericht des Wilhelm-Realgymnasiums, 1908. Val. Ernst Loescher: *M. v. Engelhart, \emph{Val. Ernst Loescher nach seinem Leben und Wirken}, ²1856. P. Schulze, \emph{Zur Neubegruendung der Dresdener Volksschule im Beginn des 18. Jahrh.} Franz Wilhelm Kockel, \emph{Aus dem Leben e. saechs. Schulmanns}, pp. 151–176, 1900. W. Mueller, \emph{V. E. Loescher} in PRE,3 1902.—Compare H. Stephan, \emph{Neuzeit} (part 4 of Krueger’s \emph{Handbuch der Kirchengeschichte}), 1909. F. Uhlhorn, \emph{Geschichte der deutsch-lutherischen Kirche} i, 1911. P. Monroe, \emph{Text-book in the History of Education}, 1905. *S. C. Parker, \emph{History of Modern Elementary Education}, 1912. F. P. Graves, \emph{History of Education in Modern Times}, 1913. F. P. Graves, \emph{A Student’s History of Education}, 1915, 21917.

One frequently hears the assertion that \emph{Spener} (1635–1705), the “father of Pietism,” inaugurated a new era in the history of religious education, and that he and the other Pietists were the first to direct the attention of the Lutheran Church to the necessity of religious instruction. Such claims,
\quellseite{156}
however, can be made only by men who are wholly ignorant of the catechetical labors of the Lutheran Church between 1530 and 1600 and between 1600 and 1676. It is true, indeed, that Spener emphatically stressed the obligation of educating the youth and energetically devoted himself to that task. But others before him had done the same; and he himself repeatedly admitted that the exemplary educational institutions of Strassburg aroused his zeal. Again it is true that he conscientiously endeavored to prompt the young to apply to heart and life what they had recognized as saving truth and thus “to bring the head down to the heart” (cf. the question “What comfort and inducement to godliness do I find in this?” with which the discussion of any topic generally concludes). But he explicitly states that the catechisms of Gotha and Quedlinburg excel by virtue of such elements of edification. It is furthermore true that he insisted on using the Bible in school and on proving every doctrine from the Bible. But the same end had been in view when, long before his time, selections of Bible passages and collections of pericopes had been prepared; the only advance made by Spener in this respect is his effort to place the New Testament into the hands of the pupils so that by looking up the passages they might become more thoroughly acquainted with the Bible. Finally, it is true that it was Spener who secured general diffusion for \emph{confirmation} as the concluding ceremony of \emph{religious instruction} and as preparation for the first communion. But (a) according to Lutheran principles the chief feature of confirmation is not what Spener stresses, the solemnity proper and the vow, but the examination of the catechumens’ faith and knowledge; this examination, however, had been reestablished in most localities during the restoration period after the Thirty Years’ War. Moreover (b), it was not Spener at all who first evaluated confirmation as the solemn conclusion of religious instruction; he found
\quellseite{157}
it in that form in a Hessian village near Frankfort-on-Main where it had maintained itself, as generally throughout Hesse, since 1539, and simply adopted it. It is clear therefore that his catechetical efforts are not to be regarded as “a voice in the wilderness” but—and this applies equally well to all other complaints and counsels of Spener—rather as the dominant note in a concert of godly suggestions and desideria. The reforms which others before him and beside him had advocated and which some had long ago enacted, were energetically sponsored by Spener and increasingly carried into effect owing to the weight of his personality; but he is not entitled to be honored as “reformer of religious instruction.” His Catechism which was based on catechetical sermons and was published in 1677 under the title “Simple Explanation of Christian Doctrine” does not even compare with some of the catechisms mentioned in § 16; it is beautiful in some respects, but its merits are neutralized by an over-abundance of dogmatical and edifying matter; as a result it exceeds in bulkiness all limits set by precedent and convenience even if we keep in mind that the book was not primarily intended for children, but for parents, schoolmasters, and pastors. His Tabulæ Catecheticæ have the same fault. Though he could, at times, warmly recommend Bible reading, he did nearly nothing for the dissemination of Biblical History.

In point of method, Spener marks a slight advance. He was, indeed, no model catechist as his own catechism abundantly shows, nor was he conversant with the leading question whereby the unknown is reached through the known. But more resolutely than all predecessors he commenced to discard the learning by rote of ready-made answers, and he introduced the “analyzing” or “dividing question” by which the subject matter was at least logically divided (cf. his Tabulæ).

\quellseite{158}
At the same time the disciples of Spener, more than he himself, started out upon dangerous paths. Confirmation was exalted at the expense of baptism. According to the Pietists, the latter constitutes only the beginning of regeneration, and confirmation is the consummation of this work of the Spirit. Thus they regarded confirmation as a renewal, on the part of God and man, of the baptismal covenant. As a result, the instruction of confirmands completely overshadowed all other forms of religious instruction; moreover, the catechist was expected to lead the catechumens through the several stages of the way of salvation in such a manner that they actually experienced conversion and regeneration and on basis of their experience vowed henceforth to be God's own. In consequence of such treatment, on the one hand, the gradual introduction into the objective acts of God was dangerously curtailed, and on the other hand there resulted an undue stressing of the subjective element of inner experience and much emotional jugglery and self-deception. Certainty of salvation was grounded, not in the objective act of God, but rather in the — inconstant, vacillating — subjective feeling of grace. Another grievous consequence was that many pietists ignored Luther's Catechism altogether in their instruction and substituted for it the socalled \emph{Ordnungen des Heils} (successive stages of the order of salvation) which were rather frequently used till 1750.

Among the catechisms based upon Spener's principles and still fairly sound in doctrine, we mention the Dresden Cross Catechism which has been in use in many places in Germany and America down to the present day, and the Flensburg Catechism. The Catechism of Herford (1690) and of Mecklenburg (1717) have kept rather free from pietistic influence. The latter, although in revised form, and here and there also the former, are used to this day. On basis of Osiander's Booklet for Communicants (1590) and J. V. Andreae's \emph{Instruction}
\quellseite{159}
\emph{for Children} (1621), Hiemer issued the Booklet for Catechumens in 1723 which, in revised form, even today forms the basis for pre-confirmation instruction in Wuerttemberg while the Brenz-Luther Catechism is employed for other religious instruction in both church and school.

In part Spener, in part also Gotha where he had been reared, and unquestionably also some opponents of dead orthodoxy influenced \emph{August Hermann Francke}, the eminent educator and theologian whose activity was a great blessing for the cause of school and religion throughout Germany. His orphans’ school soon developed into a veritable model. In accordance with the Gotha regulations, he made room for scientific branches; children, he thought, should not only write a good hand, but also be able to compose letters independently, they must know something about their natural environment, about geography and about history, boys ought to learn the elements of surveying, and boys as well as girls must be able to knit and patch. One half of the school time, however, he devoted to religious instruction; hence we meet here all the features which Ernest of Gotha had introduced in his country, also a number of improvements: more stress was laid upon training in prayer; the catechism was analyzed by questions and consistently applied to life; sections from the Bible were read and explained daily; surveys of the more important books of the Bible and of the whole Bible were given; the children were instructed in the private devotional use of the Bible; and in the Paedagogium, the higher school, Castellio’s and Fabricius’ books were used as the basis for regular instruction in Biblical History. From all parts of Germany pupils flocked to the institutions at Halle to receive their normal training in the Paedagogium or Normal School. Teachers who had been educated in Francke’s schools were found in every section of the country, and additional Normal Schools patterned after Halle were
\quellseite{160}
established in other states of Germany. Francke’s views were given widest publicity and in most cases cordially received. In fact, when they were largely incorporated in the Prussian Regulations for Schools (1763; in Saxony 1773), they had, even if with some important modifications, conquered nearly all of Germany.

In the first half of the eighteenth century instruction in Biblical History received more attention also in the common schools. School regulations prescribed it for Dresden in 1711, for all electoral Saxony in 1713, for Prussia in 1716. In 1714 Huebner, the Hamburg rector, issued “\emph{Twice 52 Stories from the Old and New Testaments}”; while this book is not, as has been claimed, the “first Biblical History”, it has been a favorite text book till far into the nineteenth century on account of the skillful selection of stories and the questions appended to them. The Prussian Regulations mentioned above fixed Saturday as the day for instruction in Bible History. Also Valentine Loescher, an eminent champion of orthodoxy, advocated the introduction of Biblical History as a branch of religious instruction, and pleaded for a general regeneration of the education of the young. In the first quarter of the eighteenth century, probably owing to Francke’s influence, the first courses in catechetical theology were given in universities with Buddeus and Rambach as lecturers.

\quellseite{161}
\section[§ 18. The Period of Rationalism]{§ 18. Religious Instruction in the Church During the Period of Rationalism (About 1750–1830)}

Ph. Schuler, pp. 226 ff. G. v. Zeschwitz ii 2, 1 ch. 17.18. *Schumann-Sperber, pp. 61–87. J. Steinbeck, pp. 48–53. K. F. A. Kahnis, \emph{Der innere Gang des deutschen Protestantismus}, ³1874. *K. v. Raumer, \emph{Geschichte der Paedagogik} ii, ⁶1897. Chr. E. Luthardt, \emph{Der Rationalismus im Gebiet der Schule}, ³1891. *M. Schian, \emph{Die Sokratik im Zeitalter der Aufklaerung}, 1900. CHRIST. WOLFF AND HIS THEOLOGY: H. Stephan, \emph{Christ. Wolff} in PRE³, 1908. JOHANN LORENZ MOSHEIM: K. Heussi, \emph{J. L. Mosheim}, 1906. JOH. BERNH. BASEDOW: G. Hahn, \emph{Basedow und sein Verhaeltnis zu Rousseau}, 1885. O. H. Lang, \emph{Basedow, His Life and Educational Work}, 1891. H. Lorentz, \emph{Entwicklung und Bedeutung der Paedagogik Basedows} in Jahn-Fleckeisens Jahrbuecher, 1893. A. Pinloche, \emph{Geschichte des Philanthropinismus}, 1895. Balike, \emph{Die Stellung der Philanthropisten zum Religionsunterricht}, 1901. J. B. Basedows \emph{Paedagogische Schriften} ed. by H. Goering, 1880. J. J. ROUSSEAU: W. Boyd, \emph{The Educational Theory of Rousseau}, 1911. His \emph{Emile}, or treatise on education is translated into English by W. H. Payne, 1893. KARL FRIEDR. BAHR T: G. Frank, \emph{K. Fr. Bahrt} in Raumers Historisches Taschenbuch, 1866. L. Leyser, \emph{K. Fr. Bahrt, der Zeitgenosse Pestalozzis, sein Verhaeltnis zum Philanthropismus und zur neueren Paedagogik}, ²1870. Moses Mendelsohn, \emph{Phaedon oder ueber die Unsterblichkeit der Seele}, 1767; also reprinted in Reclams Universalbibliothek. CHR. G. SALZMANN: *E. Ackermann, \emph{Salzmanns ausgewaehlte Schriften}, 2 vols., 21901. Besides his \emph{Krebsbuechlein} (1780, reprint in Reclam), \emph{Ameisenbuechlein} (1806, repr. in Reclam), and \emph{Konrad Kiefer} (1796), compare esp. \emph{Ueber die wirksamsten Mittel, Kindern Religion beizubringen} (1780; new edition by P. Schuetze, 1905), \emph{Die Familie Ehrenfried, oder erster Unterricht in der Sittenlehre fuer Kinder von 8–10 Jahren} (1808), \emph{Heinrich Gottschalk in seiner Familie, oder erster Unterricht fuer Kinder von 10–12 Jahren} (1804), and \emph{Christliche Hauspostille} (1792 f.). JOHANN HEINRICH PESTALOZZI: L. W. Seyffarth, \emph{J. H. Pestalozzi nach seinem Leben und seinen Schriften dargestellt}, 1872, 61903. H. Kruesi, \emph{Pestalozzi, His Life, Work and Influence}, 1875. A. Pinloche, \emph{Pestalozzi and the Foundation of the Modern Elementary}
\quellseite{162}
\emph{School}, 1901. P. Natorp, J. H. \emph{Pestalozzi}, 3 vols., 1905. W. S. Monroe, \emph{History of the Pestalozzian Movement in the United States}, 1907. J. Natorp, \emph{Pestalozzi, sein Leben und seine Ideen}, 1908, ²1912. H. Holman, Pestalozzi, 1908. J. A. Green, \emph{Life and Work of Pestalozzi}, 1913. H. Debes, \emph{Das Christentum Pestalozzis}, 1880. L. W. Seyffarth, \emph{Saemtliche Werke Pestalozzis}, 12 vols., 1899–1902. *F. Mauer, \emph{Ausgewaehlte Werke Pestalozzis}, 4 vols., ⁵1897–1906. J. A. Green, \emph{Pestalozzi’s Educational Writings}, 1912. \emph{Abendstunde eines Einsiedlers} (1780), reprinted in Ostermann, \emph{Paedagogisches Lesebuch}, 1905. \emph{Lienhard und Gertrud} (1781–1787), reprinted in Reclam; translated into English and abridged by Eva Channing, 1896. \emph{Wie Gertrud ihre Kinder lehrt} (1801) reprinted in Reclam; English translation by L. E. Holland and Fr. C. Turner, 1898. A. Israel, \emph{Pestalozzi-Bibliographie}, 3 vols., 1903. GUST. FRIEDR. DINTER: Chr. Palmer, \emph{G. Fr. Dinter} in K. A. Schmid’s Enzyklopaedie, ²1878. W. Amelung, \emph{Dinters Grundsaetze der Erziehung und des Unterrichts}, 1881. Fr. Seidel, \emph{Dinters Schriften in Auswahl}, ²1887–1889. *A. Schultz, \emph{G. Fr. Dinter, sein Leben und seine Schriften}, 1908. FR. H. CHR. SCHWARZ: K. B. Hundeshagen, F. H. Chr. Schwarz in PRE,³ 1906. F. D. E. Schleiermacher: A. Heubaum, F. D. E. Schleiermacher in Rein’s Enzyklopaedisches Handbuch der Paedagogik, ²1902. G. v. Rohden, \emph{Schleiermachers Paedagogik}, 1884. *E. Platz, \emph{Schleiermacher’s paedagogische Schriften}, ³1902.—Compare H. Stephan, \emph{Neuzeit} (part 4 of Krueger’s \emph{Handbuch der Kirchengeschichte}), 1909. F. Uhlhorn, \emph{Geschichte der deutsch-lutherischen Kirche} ii, 1911. P. Monroe, \emph{Text-book in the History of Education}, 1905. *F. P. Graves, \emph{History of Education in Modern Times}, 1913. S. C. Parker, History of Modern Elementary Education, 1912. F. P. Graves, A. Student’s History of Education, 1915, ²1917.

By yielding ever increasingly to subjectivism, the disciples of Spener prepared the soil for the rise and spread of the Wolffian School (Christian Wolff, 1679–1754) and eventually for Rationalism. Wolff and his followers attempted to apply the methods of mathematical demonstration to theology, believing in all sincerity that they would thereby fortify the doctrine of the Church against attacks. Their method, however, led them inevitably to the point where they made reason the arbiter of revelation; thus they became the forerunners of Rationalism which recognized no other source of knowledge than reason and regarded as morally good whatever
\quellseite{163}
is useful. As far as religious instruction is concerned, the earlier exponents of Rationalism left the body of Christian doctrine intact, but they deemed an analyzing process sufficient for disclosing and setting forth its vital values (Loe-secke, Baumgarten); as a matter of fact, the children were by no means always taught to think, and very little permanent influence was exerted upon their will. Later Rationalists openly broke with the doctrine of the Church, demanded the abolition of Luther’s Catechism and substituted other text books for it. This turn of affairs was in no small degree due to the prestige which the so-called philanthropist Basedow (1723–1790) had gained; he aroused much interest in education, but thought that the sole aim of education was to train useful citizens of this world, and therefore advocated the divorce of school and church, abolition of confessional instruction, and the introduction of courses in morality. His principles were cordially received far and wide. That morality is well-pleasing to God only when based upon true religion, that real power for good can grow only from a heart which is devoted to God, had been forgotten. But even in those circles where orthodox doctrine was preserved, serious defects were apparent: Luther’s Catechism was still used, indeed, but it was not organically unfolded for the children; instead they were given independent dogmatical and ethical expositions (sometimes anything but congruous) which only too frequently expunged or at least concealed the very core of Christian truth, viz., faith in the crucified and risen Saviour as the one and only root of true morality. Questions which even today haunt some catechetical expositions, such as “How many religions are there?” “What is natural religion?” “What additional sources of religion are there besides Scripture?”, and the division of “duties” into “love of God, love of self, and love of the neighbor”, are remainders of the heritage from this period. The Hanoverian Catechism of
\quellseite{164}
1790 and Seiler’s Catechism of 1779, long used in Bavaria, are among the better products of this age. Salzmann in Schnepfental (Thuringia) and Dinter in Koenigsberg (1760–1831) are strong exponents of the rationalistic type of religious instruction. Dinter at least tried to come to terms with revealed truth, but the catechetical works of Parisius, Mayer, and others contain deistic philosophy and physiology rather than Christian saving truth.

Yet, in one respect this period marks an advance beyond all the catechetical labors of the past and is to be considered as the foundation of the better conditions obtaining in the present. From what has been said, this element of progress can not be found in the sphere of the material of instruction although we meet now and then with wise and commendable restriction to the essentials of Christianity, but rather in the \emph{didactic method}. The method was that of Socrates who like the obstetrician endeavored to draw forth into the light of day the germs of truth slumbering in the minds of his disciples, and has therefore been named “Socratic Method”. Already in his Ethics of Holy Scripture Mosheim, as early as 1735, had recommended this method for religious instruction of children; but he presupposed that the matter of instruction was first imbedded in the memory; it was then to be explained through the question method with the aim of bringing about firm convictions of faith and virtuous maxims of life. Apparently, however, Mosheim’s principles did not affect the catechetical practice of the decades immediately following. The Socratic method did not win favor until the appearance of Basedow’s \emph{Method of Instruction} (1764), of the \emph{Zuerich Questions} (1772), and the \emph{Philanthropinischer Erziehungsplan} (Philanthrophy as the Basis of Education) by the frivolous and licentious K. Fr. Bahrdt (1776). It is quite likely that Moses Mendelsohn’s \emph{Phaedon} (1767) is responsible for the name given to the method. The members of the Socratic
\quellseite{165}
school, at least its earlier exponents, actually attempted by skillful questioning to draw forth from the child’s mind all articles of faith and rules of life; in their opinion it was unnecessary to put anything into the child’s mind since everything, even if somewhat hazy and inchoate, is found in it by nature. However erroneous this conception, and however inadequate and defective this method, one important result was achieved: war was declared upon all senseless, mechanical memorizing. The proposition “Nothing is to be memorized unless it is understood” was and is sometimes exaggeratively applied, it is true; but if correctly employed, it helps develop an educational method which appeals to the intellect of the child, engages his co-operation, and by proceeding from the near to the far, from the known to the unknown, succeeds in arousing his interest, i.e., his mental presence and sympathy, and in moulding his whole future development. Also in this respect Dinter’s achievements are noteworthy, especially because he, more than all other rationalists who followed the Socratic method, proceeded from Biblical foundations instead of a supposititious innate body of knowledge possessed by the child. Even today one may learn much about the developing method of teaching from his \emph{Discourses upon the Chief Parts of Luther’s Catechism} (1806 ff.).

While the Socratic method ruled undisputedly till the end of the eighteenth century, a counter movement was ushered in at the beginning of the new. Its protagonist was \emph{Pestalozzi} (\emph{How Gertrude Teaches her Children}, 1801). He had much in common with the Socratics, but he endeavored to reach not only the head but the whole man, including the heart; and dissatisfied with an occasional, more or less accidental employment of intuition (\emph{Anschauung}) he made the principle of intuition fundamental for the whole process of education. The movement was furthered by Schwarz of Heidelberg through his important \emph{Catechetics} (1818): he
\quellseite{166}
made some concessions to the Socratic method, but demonstrated clearly that after all many religious truths simply cannot be elicited from the child’s mind and that they postulate another source. The Socratic method was still further discredited when Schleiermacher (1764–1834) notwithstanding his demand for the retention of the developing method, lifted up his mighty voice against it and maintained that the acroamatic method had to go hand in hand with the developing; the teaching of general truths he declared insufficient: the children are to become Evangelical Christians who will recognize, at the end of instruction, as their own the faith of the Church summarized in the Apostolic Creed. And when finally throughout the land men returned to the faith and theology of the Bible, the Socratic period so closely linked with Rationalism was at an end.

\quellseite{167}
\section[§ 19. Catechetical Labors Since the Renewal of Faith]{§ 19. Catechetical Labors Since the Renewal of Faith. Since About 1830}

*G. v. Zezschwitz ii 2, 1, ch. 34–50; ii 2, 2, ch. 19. Schumann-Sperber, pp. 83–122. J. Steinbeck, pp. 53–63. *G. Thomasius, \emph{Das Wiedererwachen des evangelischen Lebens in Bayern}, 1867. R. Rocholl, \emph{Geschichte der evangelischen Kirche}, 1897. G. Ecke, \emph{Die evangelischen Landeskirchen Deutschlands im 19. Jahrh.}, 2 vols., 1897. 1904. H. Stephan, \emph{Die Neuzeit} (part 4 of Krueger’s \emph{Handbuch der Kirchengeschichte}), 1909. F. Uhlhorn, \emph{Geschichte der deutsch-lutherischen Kirche}, ii, 1911. F. Schindler, \emph{Kritischer Wegweiser durch die Literatur des Konfirmandenunterrichts}, 1899. F. Cohrs, \emph{Katechismusunterricht} in PRE³, 1901. H. Scherer, \emph{Fuehrer durch die Stroemungen auf dem Gebiet der Paedagogik und ihrer Hilfswissenschaften}, part ii: Religions- und Moralunterricht, 1907. J. Berndt, \emph{Methodik des Unterrichts in der evangelischen Religion}, pp. 131–140, 1909. H. Lewin, \emph{Geschichte der Entwicklung der preussischen Volksschule}, 1910. *W. Ostermann, \emph{Geschichte der Paedagogik}, pp. 120–291, ⁴1910. E. Thraendorf, \emph{Allgemeine Methodik des Religionsunterrichts}, pp. 14–36, ⁵1912. WILHELM HARNISCH: R. Rissmann, \emph{W. Harnisch in seiner Bedeutung f. d. Entwicklung der deutschen Volksschulpaedagogik}, 1887. Metzmacher, \emph{Weiter- und Umbildung der Pestalozzischen Grundsaetze durch Harnisch}, 1901. Harnisch’s \emph{Volksschulwesen} has been republished by Bartels, 1893. F. A. W. DISTERWEG: G. Voigt, \emph{Diesterweg und die evangelische Volksschule}, 1891. *E. v. Sallwuerk, \emph{Leben, Lehre und ausgewaehlte Schriften Diesterwegs}, 3 vols., 1899–1900. JOHANN FR. HERBART: W. Kinkel, \emph{Joh. Fr. Herbart, sein Leben und Wirken}, 1903. *O. Fluegel, \emph{Herbarts Leben und Lehren}, 1907. G. Voigt, \emph{Die Bedeutung der Herbartschen Paedagogik fuer die Volksschule}, ⁴1908. *Chr. de Garmo, \emph{Herbart and Herbartianism}, 1895. Chr. Ufer, \emph{Introduction to the Paedagogy of Herbart}, ed. by Chr. de Garmo, 1894. W. J. Eckhoff, \emph{Herbart’s ABC of sense-perception}, 1896. *Sallwuerk u. Bartholomaei, \emph{Herbart’s paedagogische Schriften}, 2 vols., ⁷1903, ⁷1906. H. M. and E. Felkin, \emph{Herbart’s Science of Education}, translated, with a biographical introduction, 1895. TUISKON ZILLER, T. Ziller, \emph{Grundlegung zur Lehre vom erziehenden Unterricht}, 1865. *T. Ziller, \emph{Allgemeine Paedagogik}, ²1892. *T. Ziller, \emph{Einleitung in die allgemeine Paedagogik}, ²1901. G.
\quellseite{168}
Froehlich, \emph{Die wissenschaftliche Paedagogik Herbart-Ziller-Stoys}, 1901. E. v. Sallwuerk, \emph{Das Ende der Zillerschen Schule}, 1904. Compare the short but good chapter on Herbartianism in F. P. Graves, \emph{A Student’s History of Education}, pp. 333–369, 1915, and the excellent chapter on The Herbartians in S. Ch. Parker, \emph{History of Modern Elementary Education}, pp. 375–430, 1912. F. W. Doerpfeld: A. Carnap, \emph{F. W. Doerpfeld. Aus seinem Leben und Wirken}, ²1905. E. Hindrichs, F. W. Doerpfeld, ²1906. *F. Wienstein, F. W. Doerpfeld, \emph{Sein Leben und seine Schriften}, ²1907. Doerpfelds \emph{Gesammelte Schriften}, 12 vols., ¹894 ff.; the most important writings of Doerpfeld’s are \emph{Denken und Gedaechtnis}, ¹1906, \emph{Der didaktische Materialismus}, ⁵1915, \emph{Religioeses und Religionsunterrichtliches}, ¹895, \emph{Zur Methodik des Religionsunterrichts}, ⁵1906. Compare the pertinent articles in Rein, \emph{Enzyklopaedisches Handbuch der Paedagogik}, ²1902 ff.; also W. Rein, A. Pickel, u. E. Scheller, \emph{Theorie und Praxis des Volksschulunterrichts nach Herbarts Grundsaetzen}, 8 vols., ⁵1901–1908. History and Practice of Sunday Schools in Germany: H. Dalton, \emph{Geschichte, Wesen und Weise der evangelischen Sonntagsschule}, 1887. H. von der Goltz, \emph{Das Beduerfnis besonderer Jugendgottesdienste}, 1888. Reinhard, \emph{Zur Geschichte der Kindergottesdienste und Sonntagsschulen in Deutschland}, 1888. L. Tiesmeyer, \emph{Die Praxis der Sonntagsschule}, 1877. *F. Dibelius, \emph{Der Kindergottesdienst}, 1881. *Al. Schumann, \emph{Der Kindergottesdienst in seiner gesteigerten Bedeutung fuer Gegenwart und Zukunft der evangelischen Kirche}, 1909. \emph{Der Sonntagsschulfreund}, ed. by Fleischmann, 1869 ff. Der Kindergottesdienst, ed. by Tiesmeyer, Zaule ck, and Volkmann, 1891 ff.

Amid the wreck and ruin in church and school caused by Rationalism, the old catechisms preserved their hold upon the home and supplied the pious with food. A number of \emph{catechetical} works with Biblical content had appeared during the first quarter of the nineteenth century, and with the revival of confessional Lutheranism (about 1830) genuinely evangelical principles were once more quite generally applied to the work of training the young. Men like Nitzsch, Klaus Harms, and Hengstenberg made their influence felt in Northern Germany, and Brandt, Bomhard, Thomasius, Harless, and Loehe in Southern Germany. In Prussia Harnisch combatted Dinter’s influence (Hilfsbuch fuer Lehrer, Aid for Teachers, 1834); in Saxony Rudelbach advocated the restoration
\quellseite{169}
of \emph{Christenlehre} (instruction in evangelical doctrine on Sunday afternoons; 1840); in Bavaria Ackermann defended the use of the Chief Parts of the Reformers (1832) and Boeckh issued a good exposition of the Catechism. In 1845 Loehe published his excellent Catechism; one of the chief merits of this book is that it is free from dogmatical and historical encumbrances; however pressing the temptation may have been to include such material, Loehe refused to yield to it, but rather enunciated the following important principle: An exposition of the Catechism is no more than a presentation, according to the rules of exegesis, of the treasures of religious truth contained in Luther’s Enchiridion. Jaspis in Pomerania (1850) reminded again of the fact that no better illustrative material for catechetical instruction may be found than Biblical History. Mecklenburg-Strelitz introduced a splendid catechism in 1849. Between 1850 and 1860 the number of confessional Lutheran catechisms increased steadily. Special mention must be made of Caspari in Bavaria, Ernesti in Brunswick, Seebold in Hanover, Nielsen in Schleswig-Holstein. Also within the “Union” there was renewed activity. A good catechism was issued, e.g., for Baden in 1856 which served as the basis for the Rhenish Catechism. In Wuerttemberg the Brenz-Luther Catechism and the Booklet for Catechumens were revised. The School Regulations included catechetical instruction in the curriculum, and most of the old educational institutions in church and school were resuscitated and adapted to the new conditions. Expositions of the catechisms employed in the different regions were issued for the children; and bulky volumes for teachers and pastors were written by Wangemann, Arndt, Kaehler, Nissen, and others. Even catechetical sermons were preached once more; those of Loehe, Heubner, Ahlfeld, and Caspari are particularly excellent. The origin of Luther’s Catechism was investigated, the earliest attainable editions were republished
\quellseite{170}
(Harnack, Moenckeberg, Schneider); and collections of stories were issued as illustrative material (e.g., Caspari, \emph{Geistliches und Weltliches}). Of the expositions issued between 1860 and 1870 Buchrucker’s Explanation and his handbook \emph{Der Katechismusunterricht} in which he strongly pleads for the developing method, are the best. In the subsequent decades an almost endless variety of catechetical works appeared. The best books for children are those by Steinmetz and Beck, for the use of pastors those by Schuetze, Muenchmeyer, Schumann, Zezschwitz, and Steinmetz; Th. Kaftan’s exposition, however, deserves the palm for its discriminating and exquisite presentation of material. The most instructive product of the Ritschlian school is B. Doerries, Der Glaube. At present, certain theologians and pedagogues again demand with loud voices the elimination of Luther’s Small Catechism from the instruction of children. Compare § 35.

Another result of the renewal of faith was the increasing appreciation of \emph{Biblical History} as educational material. In their conflicts with the Rationalists the champions of the old faith stressed the fact that Christianity is based upon divine revelation, and quite naturally they came to evaluate the introduction into the history of revelation, into the genesis of Christianity, as an integral part of the education of the young. Especially Harnisch and F. L. Zahn advocated this branch. Harnisch emphatically declared that instruction in the Catechism can be profitably given only where Bible History has first been taught. In the common schools a number of stories were to be devotionally explained and recited; the teachers were given many valuable hints in Nissen’s Unterredungen ueber die biblische Geschichte (Talks upon Biblical History). Huebner’s stories of 1714 were widely used as text-book, also those issued by the Calw Publishing House. In the higher schools Biblical History was expanded into a history of the kingdom of God; here the books by Harnisch,
\quellseite{171}
Zahn, and later by Kurtz and Thomasius were often used. Franz Wiedemann’s How I Tell My Children the Stories of the Bible assisted mothers and other teachers of the little ones. In proportion as the views of Hofmann, the great Bible theologian of Erlangen, spread, it was deemed desirable even in the common schools, at least in the higher classes, to connect the individual stories of the Bible into a connected history of salvation. It was particularly Buchrucker who advocated this plan and composed excellent manuals for this purpose, a Biblical History for the pupils, and \emph{Biblischer Geschichtsunterricht} for the teacher. Where this seemed unattainable, the upper classes would at least study some outstanding Bible characters, as Wangemann had advised. At present Biblical History is so strongly emphasized in some quarters that all instruction in the Catechism is to be excluded until the pupils reach the higher classes. Of late “modern” theology and pedagogy begin to affect also instruction in Biblical History: the value of the Old Testament and in part the New Testament as records of the history of redemption is denied, yes even their historicity is frequently rejected; large parts, such as the account of creation or of the patriarchs, even the narratives of Jesus’ infancy and of his resurrection, are dismissed as legendary; there are many who expressly demand fairy tales as material of religious instruction in the lower grades and others desire the child to experience inwardly all the stages of religious “evolution” from naïve paganism to evangelical Christianity (Kulturstufentheorie). Compare § 34.

\emph{Instruction in the Bible} was also retained in this period, both in the form of Collections of Bible Passages, used rather exclusively as proof-texts for the Catechism, and in the form of Bible reading and of Bible Literature. A question that was much discussed was whether or not children should be given a complete Bible, or a School Bible or Biblical Reader
\quellseite{172}
which might be replaced by a complete Bible at the close of the school period. The \emph{Biblical Reader} by Voelker-Strack, and those for Bremen and for Wuerttemberg have found wide circulation.

Also the \emph{Hymnal} continues to be used in instruction. Well known hymn-tunes are practiced, and owing to the efforts of Lauritz and J. Zahn, in their original rhythmic setting; a treasure of choice hymns (the reduction of which to a smaller number is frequently demanded) is learned by heart and explained; the liturgy of the chief service is discussed, and often the pericopes are read and briefly interpreted. Also the most important events of \emph{Church History} are now usually treated in school, thus realizing one of Augustine’s ideals. Increasingly the need is felt of instruction concerning the polity and constitution of the Church, her tasks and duties in the present time, the work of missions, etc.; but specific forms for such instruction have not as yet been developed (cf. Warneck, \emph{Missions in the School}). The conviction is gaining ground that one cannot do justice to all these topics in the time previous to confirmation. They are therefore often studied in \emph{Christenlehre} where this ancient institution still exists, also in the Young People’s or Young Men’s Societies which owe their origin to the Church’s interest in the young after their confirmation. Others would remedy the situation by changing the confirmation practise or at least the time of confirmation so that confirmation will not any more coincide with graduation from grammar school. The supplementary care of the confirmed is still largely an un-solved problem which, like many other problems in the sphere of religious instruction, is highly complicated by the recent recrudescence of unbelief.

Since 1863 \emph{Sunday Schools}, introduced from America, have been built up in many cities of Germany. In most cases
\quellseite{173}
they have developed into religious \emph{services for children}, with or without division into groups.

Much was accomplished during the last century in the field of \emph{method}. The fundamentals of the Socratic Method were retained, but the attempt to elicit all truths from the child was made no longer; instead the material was first presented to the child, then made clear by the developing method, and thus utilized as an incentive for his will and food for his heart. During the whole century predominant influence was exerted by the great pedagogue Herbart (1776–1841; \emph{Umriss paedagogischer Vorlesungen}, 1835) who had been inspired by Pestalozzi. Owing to his and his followers’ efforts the peculiarities of the child’s soul life were given due — often one-sided — attention, as is seen from the \emph{Formalstufentheorie} (compare § 32). The progress made during the last century in the education of the German public school teachers rendered possible a fairly universal application of this method to religious instruction. But the theologically trained catechists rather avoided it; they tried to reach their aim of gripping and permanently influencing the whole soul by a combination of the acroamatic and erotematic methods and by adapting the educational material, as the catechism, and hymns, and esp. Bible History, to their pupils’ mental ability. Among the outstanding educators it was chiefly Doerpfeld who reduced the Herbartian theory to sane limits.

Of works embracing the whole sphere of Catechetics we mention the text-books by Kraussold (1843), Palmer (1844), Zezschwitz (1863–1872), Schuetze (1876), Kuebel (1876), Th. Harnack (1882), Buchrucker (1889), Sachsse (1897), v. Nathusius (1904), Gottschick (1908), and Steinbeck (1914), to which must be added the sections on Catechetics in the textbooks of Practical Theology, especially those of Krauss, Knoke, and Achelis. O. Baumgarten’s Neue Bahnen (1903) is suggestive and instructive; like Kabisch’s \emph{Wie lehrt man}
\quellseite{174}
\emph{Religion?} It discusses the subject from the standpoint of modern theology. The works on Didactics or Public School Methods by Doerpfeld, Kehr, Schumann, Staude, Thraendorf, Rein, Reukauf, and others contain much valuable material on methods of religious instruction. For the history of catechetical methods compare also § 34–36.

\quellseite{175}
\section{§ 20. Catechetical Labors in America}

G. Weber, \emph{Geschichte der akatholischen Kirchen und Sekten in Grossbritannien}, 1845f. \emph{Original Letters relative to the English Reformation}, edited by the Parker Society, vol. 2, 1848. \emph{Three Primers put forth in the Reign of Henry VIII.}, 1848. \emph{First Prayer Book of Edward VI compared with successive revisions}, 1883. A. F. Mitchell, \emph{Catechisms of the Second Reformation}, 1886. H. Bonar, \emph{Catechisms of the Scotish Reformation}, 1886. *B. M. Schmucker, \emph{Luther’s Small Catechism, The Editions, Translations, and Explanations Published in America} in Lutheran Church Review 1880. *H. E. Jacobs, \emph{The Lutheran Movement in England}, 1890. P. L. Ford, \emph{The New England Primer}, reprint of the earliest known edition, 1899. E. F. K. Mueller, \emph{Bekenntnisschriften der reformierten Kirche}, 1903. M. Reu, \emph{Quellen zur Geschichte des kirchlichen Unterrichts}, 1 and 2, 1904, 1911. A. Lang, \emph{Der Heidelberger Katechismus und verwandte Katechismen}, 1907. F. Watson, \emph{English Grammar Schools before 1600}, 1908. W. H. Free and W. H. Kennedy, \emph{Visitation Articles and Injunctions of the Period of the Reformation}, 1910. John W. Prince, \emph{Wesley on Religious Education}, 1926. L. G. Pray, \emph{History of Sunday Schools and of Religious Education}, 1847. I. C. Power, \emph{The Rise and Progress of Sunday Schools, a Biography of Rob. Raikes and William Fox}, 1863. S. Gilbert, \emph{The Lesson System: the Story of its Origin and Inauguration}, 1879. H. C. Trumbull, \emph{Yale Lectures on the S. S.}, 1888. J. H. Harris, \emph{Rob. Raikes, The Man and his Work}, 1899. M. C. Brown, \emph{Sunday School Movement in America}, 1901. S. B. Haslett, \emph{Pedagogical Bible School}, 1903. O. G. Michael, \emph{Sunday School in the Development of the American Church}, 1904. F. W. Rice, \emph{Origin and Expansion of the S. S.}, 1906. *H. F. Cope, \emph{The Modern S. S. in Principle and Practice}, 1909. *H. H. Meyer, \emph{The Graded S. S. in Principle and Practice}, 1910. H. F. Cope, \emph{Evolution of the S. S.}, 1911. J. R. Sampey, \emph{International Lesson System}, 1911. J. R. Hurlbut, \emph{Organizing and Building up the S. S.}, 1911. R. P. Shepherd, \emph{Religious Pedagogy in the Modern S. S.}, 1911. Proceedings of the Religious Education Association 1903 ff. \emph{The Encyclopedia of Sunday Schools and Religious Education}, 3 vols., 1915. De Garmo, \emph{Principles of Religious Education}, 1901. *G. U. Wenner, \emph{Religious Education and the Public School}, 1907. A. N. Riley et al., \emph{The Religious Question in}
\quellseite{176}
\emph{Public Education}, 1911. H. Chapell, \emph{The Church Vacation School}, 1915. *W. S. Athearn, \emph{Character Building in a Democracy}, 1924. *F. O. Gift, \emph{Week Day Religious Education}, 1926. State Teachers College of Colorado, \emph{Special Bulletin of Religious and Moral Education}, 1913. *S. C. Parker, \emph{History of Modern Elementary Education}, 1913. W. P. Burris, \emph{The Public School System of Gary, Ind.}, U. S. Bureau of Education Bulletin 591, 1914. W. J. Mutch, \emph{Religious Day School} in Encyclopedia of S. S. etc., 1915. R. S. Bovine, \emph{Daily Vacation Bible School Association}, ibid., 1915. *H. H. Meyer, \emph{Cooperation in Christian Education}, 1917. S. A. John, \emph{Die Sonntagsschule, ihre Bestimmung u. Arbeit}, 1897. C. W. Hertzler, \emph{Die religioes-sittliche Erziehung d. kirchl. Jugend}, 1913. F. W. Schneider, \emph{Barclays Handbuch zur Heranbildung von Sonntagsschullehrern}, 1916. H. E. Jacobs, \emph{History of the Ev.-Luth. Church in the U. S.}, 1893. A. Graebner, \emph{Geschichte d. luth. Kirche Amerikas}, 1892. G. I. Fritschel, \emph{Geschichte der luth. Kirche Amerikas}, 1896f. J. L. Neve, \emph{Kurzgefasste \emph{Geschichte der luth. Kirche Amerikas}}, 1915 (here further literature). J. C. W. Lindemann, \emph{Schulpraxis}. J. Schaller, \emph{Notwendigkeit d. christl. Gemeindeschule fuer die christl. Familie, Kirche und Staat}. G. M. Grossmann, \emph{Die christl. Gemeindeschule}, 1895. J. R. E. Hunt, \emph{The Luth. S. S. Handbook}, 1911. R. Bunge, \emph{Weide meine Laemmer}, 1912. E. H. Engelbrecht, \emph{Why Lutheran Parish Schools}, 1915. G. H. Gerberding, \emph{The Lutheran Catechist}, 1910. J. H. Hertzer, \emph{Evang.-luth. \emph{Katechetik}}, 1911. M. Reu, \emph{Katechetik}, 1915. *\emph{Symposium on the Sunday School} in Lutheran Church Review, 1896. Theo. Schmauk, \emph{The General Council Sunday School System of Grading and Instruction}, 1897 ff. M. Reu, \emph{Wesen u. Aufgabe der Sonntagsschule} in Kirchl. Zeitschrift, 1900. 1901. M. Reu, \emph{Grundsaetze zur Herstellung von Sonntagsschulliteratur}, ibid., 1911. M. Reu, \emph{Wartburg Lehrmittel fuer ev.-luth. Sonntagsschulen}, 1911 ff. M. Reu, \emph{Wartburg Lesson Helps for Luth. Sunday Schools}, 1914 ff. Chr. P. Wiles, \emph{The Challenge of the S. S.}, 1916. *M. Reu, \emph{A New English Translation of Luther’s Small Catechism}, 1926.

The same inconclusiveness which characterizes the religious history of England between the years 1530 and 1550 confronts us also when we investigate the status of religious instruction in that period. Marshall’s \emph{Goodly Primer} of 1534 was uncompromisingly evangelical. Also the second edition (1535) which contains unabbreviated Luther’s \emph{Short Form} of 1520, is permeated by the \emph{sola gratia}; it is true that in
\quellseite{177}
the Litany which bears the superscription “\emph{\emph{The Invocation of the Saints}}”, a concession is made to the weak, but in the preface (pp. 123 ff.) it is explicitly stated that it is a concession, no more. Romanizing tendencies are not in evidence; the enumeration of the prohibition of image worship as a specific commandment (Dialogue, pp. 216-221) rather indicates Swiss influence. Henry VIII’s injunction of 1536, however, in no way transcends those of medieval times since the necessity of explaining the parts is ignored; the injunction reads as follows: “The curates shall, in their sermons, deliberately and plainly recite of the said \emph{Pater Noster}, the Articles of our Faith, and the Ten Commandments, one clause or article one day, and another another day, till the whole be taught and learned by little; and shall deliver the same in writing, or show where printed books containing the same be to be sold to them that can read or will desire the same.” In 1537 Archbishop Cranmer and a commission of bishops edited an authorized explanation of the Lord’s Prayer, the Ave Maria, the Creed, the Decalogue, and the seven Sacraments under the title \emph{The Godly and Pious Instructions of a Christian Man} (commonly called Bishop’s Book), but the explanation was Catholic, and not Evangelical. Of the same type is \emph{The Necessary Doctrine and Condition of a Christian Man} (also called King’s Book), an explanation issued in 1543. When Henry VIII died, in January 1547, and Edward VI acceded to the throne, more resolute measures of reformation were taken. In 1548 Cranmer issued an English translation of Justus Jonas’ Latin version of the Nuremberg Sermons for Children (Osiander whose niece had married Cranmer in 1532, was their chief author) under the title \emph{A Short Instruction into the Christian Religion; for the Singular Commoditie and Profite of Children and Young People}. Thus Luther’s Catechism and its best extant exposition was made accessible to the English people no less than to most
\quellseite{178}
of the other European peoples (see Reu, Quellen, etc.). It was not Luther’s Catechism, however, that gained vogue in subsequent years, but the \emph{Instruction to be Learned of Every Child}, a catechism which had been issued as part of the Order of Confirmation appearing in the first edition of \emph{the Book of Common Prayer} (1549), and which had to be recited from memory at the time of confirmation. It begins with a statement concerning sponsorship and then succinctly explains Creed, Decalogue, and Lord’s Prayer. After 1552, owing to Reformed influence (Butzer had been brought into the country, Calvin had become the chief adviser of Edward and Cranmer; but cf. also Marshall’s Primer), the commandments which hitherto had been compressed into short sentences, were given in the exact wording of Exod. 20. After Bishop Overall (1604) had added the section pertaining to the Sacraments, this catechism remained unaltered, and it is in force in the Anglican (Episcopal) Church to this day; it is, however, no longer printed as part of the Confirmation service, but as an independent section of \emph{the Book of Common Prayer} (since 1662). Authorized by the king, Poinet issued a larger catechism for adults and teachers in 1553, but it never gained wide circulation. Greater influence was exerted by the catechism issued by Nowell with the approval of archbishops and bishops, but it likewise failed to become an official text book for instruction; it is noteworthy, however, that Nowell’s catechism furnished the basis for the above-mentioned section on the Sacraments added to the Catechism in 1604. For many years, if we may judge by the curricula extant, Nowell’s catechism in its original Latin form as well as in English and Greek translation was used for instruction in religion and in the ancient languages.

In Scotland it was Knox who emphatically advocated religious instruction of the young. In his Book of Discipline (1560) we find the following rule for all places not large
\quellseite{179}
enough to support day schools: “Either the reader or the minister must take care of the children and youth of the parish, instructing them in their first rudiments, and especially in the catechism, as we have it now translated in the Book of our Common Order, called the \emph{\emph{Order of Geneva}}.” He also enjoined that on Sundays “after noon the young children must be publicly examined in their catechism in audience of the people, and in doing this the minister must take great diligence, to cause the people to understand the questions propounded, as well as the answers, and the doctrine that may be collected thereof.” Nor did he forget home instruction; he wrote: “Every master of household must be commanded either to instruct, or else cause to be instructed his children, servants and family, in the principles of the Christian religion.” The catechism to which he referred is Calvin’s Catechism of 1545 (cf. p. 122). The Presbyterian Church of England, subsequently also the Church of Scotland introduced the Shorter Westminster Catechism sanctioned by parliament in 1648 which is still in use. While the Large Westminster Catechism (also approved by parliament in 1648) resembles a textbook of dogmatics rather than a catechism (it is based on the \emph{\emph{Compendium theologiae}} of the Basle theologian Wolleb, issued in 1628 and contains 196 questions and answers in part very lengthy), the Shorter Westminster Catechism has justly been praised for its clearness and careful wording and its commendable restriction to essentials, but with its 107 questions it is, like the Heidelberg Catechism, still too voluminous and doctrinal. It contains a large number of proof-texts. The title of the authorized edition is: “The Shorter Catechism, agreed upon by the Assembly of Divines at Westminster, with the assistance of commissioners from the Church of Scotland, as a part of the covenanted uniformity in religion betwixt the churches of Christ in the kingdom of Scotland, England and Ireland, and approved
\quellseite{180}
anno 1648 by the General Assembly of the Church of Scotland, to be a directory for catechizing such as are of weaker capacity, with the proofs from the Scripture.” It is divided into two parts, (1) What man is to believe concerning God, and (2) What duty God requires of man. The first part is joined rather loosely to the Creed, the second part discusses the Ten Commandments, faith and repentance, the Sacraments, and the Lord’s Prayer; at the end of the book the words of the Creed, the Decalogue, and the Lord’s Prayer are given. — The Westminster Catechism was used for a long time also by the Congregationalists. — In Holland and in the Palatinate the Reformed Church used the Heidelberg Catechism which at Dort (1618) had been made the confession of all Reformed churches.

From these premises conclusions are readily drawn regarding the catechetical helps and the manner of religious instruction obtaining in America among the immigrants coming from England, Holland, and the Palatinate. The families that had been members of the Episcopal Church brought with them the Book of Common Prayer and instructed their children prior to confirmation in the Catechism of 1549 which was contained in it. Presbyterians coming from Scotland, Ireland, or England, also the Congregationalists brought the Shorter Westminster Catechism. The Dutch Reformed and the Palatines brought the Heidelberg Catechism. It is likely that the largest part of instruction was given in the home, but in accordance with the usage in the mother countries, also Sunday was devoted to religious instruction. We hear, for instance, that at Roxbury, Mass., the young were instructed on Sunday as early as 1674, at Norwich, Conn., in 1676, and in the Pilgrim’s Church, Plymouth, Mass., in 1680. At Plymouth it was resolved “that the deacons of the church be requested to assist the minister in teaching the children during the intermission on the Sabbath” (i.e., between the
\quellseite{181}
forenoon and afternoon services). The Rev. M. Jones opened a Sunday School at Newton, L. I., N. Y., in 1683. The Schwenkfelders in Berks and Montgomery counties, Pa., had a similar institution in 1734. It is claimed that John Wesley during his stay in America (1735-1738) established a Sunday School at Savannah, Ga. The Church of the Brethren at Germantown, Pa., gave its youth instruction on Sundays as early as 1738; it is said that Christopher Saur printed tickets or cards for use in this school. In 1740 Rev. Joseph Bellamy opened a Sunday School at Bethlehem, Conn., which exists to this day; his Memoir states that “he is believed to have been the first pastor in the land, if not in the world, who began and through all his ministry kept up, a Sabbath school in his congregation, regularly spending an hour in the interval of public worship, on the Sabbath, in catechizing and instructing one class of children and another of adults, in the word of God”; this is indeed an overstatement, but it furnishes proof that at that time such instruction on Sundays was not yet widely prevalent. Ludwig Hoecker, a Seventh Day Baptist or Dunker, established a Sunday School at Ephrata, Lancaster county, Pa., in 1740 in order “to give instruction to indigent children as well as to give religious instruction to those of better circumstances”. In 1744 Mrs. Greening opened a school in Philadelphia. Among the records of the early Presbyterian church at Westerly, R. I., there is a note dated 1752 according to which it was required “statedly to hear the children read a portion of ye Holy scriptures, and repeat ye Assembly’s Catechism”.

Besides the Catechisms enumerated above, of which the Shorter Westminster Catechism was specifically mentioned in the last quoted note, and the Bible, the catechetical helps most frequently used in home and school were the writings of Cotton, Harris, and Watts. John Cotton arrived in America in 1632 and published a catechism in 1646 which bore the
\quellseite{182}
title “\emph{Milk for Babes, Drawn out of the Breast of both Testaments, chiefly for the spiritual nourishment of Boston babes in either England, but may be of like use for any children, a Catechism}.” This little writing was reprinted even in England (Cambridge 1656), and it went through very many editions in New England. The author’s grand-son called it “peculiarly the catechism of New England” and said: “It will be valued and studied and improved, until New England ceases to be New England.” It does not surprise, therefore, that we find it incorporated in the \emph{\emph{New England Primer}}, the book which for over a century was the most important school book of New England. The \emph{\emph{New England Primer}} was composed probably by Benjamin Harris, and very likely it was printed to take the place of his \emph{The Protestant Tutor}. The first edition must have appeared between 1687 and 1690 as the “second Impression of The \emph{\emph{New England Primer}}, Enlarged” was advertised in 1691. The contents of the book vary in different editions, but generally it contained the simple alphabet, also rhymed and illustrated alphabets, the vowels, consonants, double letters, italic and capital letters, “easy syllables for children,” forming a syllabary of words from one to six syllables in length, the Lord’s Prayer, the Apostles’ Creed, pictures of birds, animals, and fish, each with a rhyme, “lessons for youth,” which are admonitory Scripture verses, Dr. Watts’ cradle hymn, “verses for little children” and prayers and advice for them, the Westminster Assembly’s Shorter Catechism, John Cotton’s “Spiritual Milk for Babes,” and a “Dialogue between Christ, Youth, and the Devil.” Some editions included proper names of men and women, the Ten Commandments, and the names of the books of the Bible given in their order. This primer reminds one strongly of school editions of Luther’s Catechism from the period when it was the sole school book in Germany.

Among the contributors to the \emph{\emph{New England Primer}}
\quellseite{183}
we meet also the third author who powerfully influenced school and home through his writings, the well known hymn writer Isaac Watts of whom Samuel Johnson said: “For children, he condescended to lay aside the scholar, the philosopher, and the wit, and to write little poems of devotion, and systems of instruction adapted to their minds and capacities, from the dawn of years through its gradations of advance in its morning of life.” His juvenile hymns published 1715 with the title \emph{Divine Songs} and republished 1720 with additions under the title \emph{Divine and Moral Songs for the Use of Children} had profound religious effects. Owing to their child-like simplicity and their spiritual warmth they enjoyed phenomenal circulation. And not only in England but also in America they constituted a most important factor in religious education of the young. More than 100 editions had been issued by 1750; Wilmer (\emph{Life of Watts}, p. 372) states that at the end of the century about 80,000 to 100,000 copies were printed annually; and Julian (\emph{Dictionary of Hymnology}, first edition, 1892) calls this hymn-book and two other writings of Watts which will be mentioned directly “the most popular text-books for religious education 50 years ago.” The two other important works of Watts are his \emph{Plain and Easy Catechisms} published in 1730 and his \emph{Scripture History} of 1732. The “First Catechism” presents in brief form the Christian way of salvation and a “Catechism of Scriptural names” (Who was Adam? Who was Eve? etc.); the “Second Catechism” contains an exposition of the Ten Commandments and of the Sacraments and some prayers, among them the Lord’s Prayer. His \emph{Scripture History} undoubtedly is the first Bible History in English; his collection of hymns is indeed not the first juvenile hymnal — Spangenberg had issued one as early as 1544, and several other hymnals for children had appeared after that, — but the first book of the kind to be published in the English language.

\quellseite{184}
New impulses for the religious education of the young proceeded from the \emph{Sunday School Movement} called into being by Robert Raikes of Gloucester (1735 or 1736-1811) and soon afterward transplanted to America.

\begin{quote}
E. W. Rice, in \emph{Schaff-Herzog Encyclopedia} vol. xi, summarizes the results of Harris’ investigations relative to Raikes’ life and work as follows: “Robert Raikes” (Sept. 14, 1735 is the date generally given for his birth; in the baptismal records of the Church of St. Mary de Crypt at Gloucester, however, the following entry is made: “Sept. 24, 1836, Robert, son of Robert and Mary Raikes, of this parish”; is it possible that he was a year old before he was baptized?) “was editor and proprietor of the \emph{Gloucester Journal}, Gloucester, England, and a strange mixture of the “dandy” and of the Reformer. Before he was of age, he began to visit the two prisons of Gloucester, to relieve the horrors of prison life, and to reform the prisoners. His sympathies were widened, his charities deepened; his failures in prison reform set him thinking, until he reached the conclusion that ‘vice is preventable.’ Twenty-five years later, when he was forty-four years of age, he began a ‘new experiment’ as he called it, of ‘botanizing in human nature.’ Going into the suburbs of the city, where many youths were employed in the factories, his heart was touched by the groups of ragged, wretched, cursing children. He knew their parents, homes, and habits; none ever entered the House of God. It was useless to appeal to such parents. He had tried to reform adults and had failed. George Whitefield had tried to reach the masses in Gloucester, but with meager results. Raikes was moved, therefore, to apply his maxim that ‘vice is preventable.’ ‘Begin with the child, for idleness is the parent of vice,’ and ‘ignorance is the cause of idleness’; therefore, ‘begin by instructing the child.’ These seem like trite statements now, but they were the result of long, deep thought by Raikes. Here was his mission. But the Rev. Thomas Stock, head-master of the Cathedral school, whom he met in his walk, was the man to help him. He started his first Sunday School in Sooty Alley in 1780, paying Mrs. Meredith for teaching the wretched little street children whom he persuaded to come to her kitchen for instruction. Mrs. Meredith found the boys ‘terribly bad,’ and soon the pupils were transferred to Mrs. Mary Critchley’s home in Southgate Street, whose house extended to Grey Friars. Bad as the boys were, the ‘girls were worse.’ The children were required to come with clean hands and faces, hair combed, and
\quellseite{185}
with such clothing as they had, though shoes and clothes were sometimes provided. The boys were ‘strapped’ or ‘caned’ by Raikes himself, for misbehavior; the girls were subdued by other means. The children were to remain in school from ten to twelve, then go home; to return at one, and, after a lesson, to be conducted to church; after church service to repeat portions of the catechism; then to go quietly home about five, without playing in the streets. Attentive scholars received rewards of Bibles, Testaments, books, combs, shoes, and clothing. The head teachers were paid a shilling a day. Raikes engaged four women in his schools, and procured other employment for them as rewards of diligence which ‘may make it worth sixpence more.’ The Rev. Thomas Stock ‘went around to the schools Sunday afternoon,’ says Raikes, ‘to examine the progress made, and to enforce order and decorum among such a set of little heathen.’ The boys were in classes of five, the advanced pupils acting as ‘monitors’ or teachers, teaching the younger pupils their letters. The girls in a separate room, with white tippets on their shoulders and white caps on their heads, were in classes also, with ‘monitors’ or sub-teachers over them. The monitors and sub-teachers were unpaid and voluntary, selected and directed by the paid master or mistress. For about three years Raikes looked upon his schools as an experiment. When William Fox, William Wilberforce and the Wesleys—John and Charles—and Jonas Hanway came as the guests of Raikes and his neighbors, he explained his plan, asked their counsel, and took them to the school to hear the children repeat prayers, the catechism, answer Bible questions, and sing Watts’ hymns. It is recorded that they were astonished, ‘caught the fire,’ and extended the movement. When satisfied that his scheme had passed the experimental stage, Raikes published a brief notice of it in the \emph{Gloucester Journal} of Nov. 3, 1783, which was copied into the London papers. \emph{The Gentleman’s Magazine} also published a letter of Raikes of November 25, 1783, in full, and a little later, another description by Raikes was given in the \emph{Arminian Magazine} edited by John Wesley. These and many other published accounts extended knowledge concerning the new movement, while many pulpits repeated the story and praised the institution.”
\end{quote}

Raikes had no intention to build up schools for all children, but he wanted to lift the neglected, wretched, unrestrained children of the city out of their misery; and although he was not chiefly concerned with religious instruction
\quellseite{186}
, but their mental and moral training, he gladly employed such aids as Bible reading, Catechism study, singing of hymns, and attendance of divine services. Moreover, it is very probable that in the course of time other children beside those whom he had in view originally, began to attend his schools because at that time England had no state-supported free public schools (until 1833 elementary education in England was left to the home; the Church, or private schools charged tuition). Certain elements in the Church at first violently resisted the movement: the archbishop of Canterbury called a convocation to determine methods of stopping the movement; William Penn seriously considered introducing a bill in parliament “for the suppression of Sunday Schools”; and in Scotland laymen who taught in Sunday Schools were held guilty of “Sabbath breaking.” But the movement was not to be retarded: in a surprisingly brief span of time the thoughts of Raikes had swept over the country, opposition was brushed aside, and schools began to spring up in large numbers. Bolton had a Sunday School with 80 volunteer teachers as early as 1787, the Sunday School at Stockport had 30 teachers in 1794, and in 1800 there were 156,400 Sunday School pupils.

William Fox, a Baptist, vigorously promoted the Sunday School cause and called a representative meeting in London on Sept. 7, 1785; its chairman was the philanthropist Jonas Hanway, and it was here that “\emph{The Sunday School Society}” was founded “for the support and encouragement of Sunday Schools in the different counties of England.” In 1803 “\emph{The British Sunday School Union}” was organized in London. Though founded originally by a group of Sunday School teachers for mutual assistance and encouragement, it soon began to consider as its chief purpose the establishment of new schools. Almost from its beginning this society has been a publishing agency, issuing lesson plans, lesson helps, etc.;
\quellseite{187}
it also arranged lecture courses. In 1818 the official figures for England were 5463 Sunday Schools with 477,225 pupils. Various factors contributed to this extraordinary development; H. F. Cope mentions the following four: (1) The awakening interest in the general education of the young which arose in various places; (2) a wide-spread development of humanitarian sentiment which led to the organization of many important relief and betterment societies; (3) the remarkable religious revival which is today best remembered in the work of Whitefield and Wesley; (4) the industrial revolution, together with upheavals of the French revolution and American independence.” The Sunday Schools of England have continued on the general pattern of the schools of Raikes; Cope enumerates the following characteristics: (1) Organized and conducted independently of church control; (2) without denominational oversight and promotion; (3) designed to combine elementary general education with religious instruction; (4) lacking the urgency of a secular system of general education.

The waves of the new movement soon touched America. Yes, here the Sunday School developed much more vigorously and rapidly than in England, and from here powerful impulses have been communicated to other countries. The stupendous growth was caused, (a) on the one hand by the fact that in the United States the government at first paid very little attention to the training of the young, leaving elementary education to the family and the Church. It was estimated that in Indiana less than one-sixth of the children attended any school in the first quarter of the nineteenth century. In the larger cities the Lancasterian system of monitors was applied; in Philadelphia, e.g., in 1834 there averaged one teacher to 218 pupils, the work being directed by the teacher with the aid of “monitors” who were responsible for certain groups of pupils. Such herding of pupils in large numbers was, of course, extremely ineffectual.
\quellseite{188}
The government of Massachusetts required towns to maintain schools by the eminently important law of 1647. This law reads, in part, as follows: “It being one chief point of that old deluder, Satan, to keep men from the knowledge of the Scriptures, as in former times, by keeping them in an unknown tongue, so in these latter times, by persuading from the use of tongues, that so at last the true sense and meaning of the original might be clouded by false glosses of saint-seeming deceivers, that learning might not be buried in the grave of our fathers in church and commonwealth, the Lord assisting our endeavors—It is therefore ordered that every township in this jurisdiction, after the Lord has increased them to the number of fifty householders, shall then forthwith appoint one within their town to teach all such children as shall resort to him to write and read, whose wages shall be paid either by the parents or masters of such children, or by the inhabitants in general, by way of supply, as the major part of those that order the prudentials of the town shall appoint; providing those that send their children be not oppressed by paying much more than they can have them taught for in other towns.” Parker shows that many of the Massachusetts towns failed to obey the spirit of the law of 1647. The reading texts used were religious books like the \emph{New England Primer} and similar primers, until spelling books, above others \emph{Webster’s Speller} of 1783, took their place.—“In colonial Pennsylvania, in contrast with Massachusetts,” Parker says (pp. 62 f.), “no general system of public schools was developed, but elementary education remained entirely in the hands of the churches and neighborhood organizations which were actuated by religious motives. Wm. Penn had contemplated the organization of a system of public schools, but his Utopian ideals were not realized. The second general assembly of the colony (1683) passed a law requiring that all children be taught so that they could read the Scriptures and write by twelve years of age. The law soon became a dead letter, however, owing to changes in the government and the conflicting interests of such a cosmopolitan colony. The assurance of liberty of religious worship attracted in Pennsylvania a great many Protestant religious emigrants and exiles. These included Quakers, Baptists, Presbyterians, Methodists, German Lutherans, members of the German Reformed Church, Moravians, and others. They were generally believers in the Protestant principle of training to read the Bible as the road to salvation, and each group of religious enthusiasts set up a school as an essential part of its religious organization. These church schools predominated in the
\quellseite{189}
eastern portions of the state. In the more thinly settled and frontier parts of the state, where these compact religious communities were not found, the more mixed communities tended to establish subscription (voluntary), or ‘neighborhood’ schools. These were generally the result of the voluntary cooperation of a few families, often stimulated by some energetic and wealthy father who desired that his children should have at least an elementary education. The neighborhood schools were most common in the western part of the state. Together with the Church schools they provided nearly all the elementary education available down to 1834.”

On the other hand, (b) the success of the American Sunday School was due to the fact that the Church fostered the new institution, regarding it from the very beginning as the real church school and subsequently as the one and only church school. The earliest American Sunday Schools were quite similar to those founded by Raikes; the “First Day or Sunday School Society” founded in Philadelphia in 1791 by Bishop White (Episcopalian), Mathew Carey (Roman Catholic), and Benjamin Rush (Universalist) was organized in order to maintain schools for those “who on Sundays were employed in the worst of purposes, the depravity of morals and manners,” but soon the sphere of activity was widened so as to embrace all children. Again, wherever no elementary public school existed or where the Church had not established any day school, the Sunday School also gave some elementary instruction, but this feature disappeared when the number of public state-supported schools began to increase, and the Sunday School became a school exclusively religious. And when the ever improving secular school had crowded out of existence the religious day school the Sunday School became the one and only school of the Church because religious instruction was excluded from the public school. Small wonder that many hands were busily engaged in improving, extending, and systematizing the work of the Sunday School.

It appears that the Methodist Church was the first to cultivate the Sunday School in its newer form in the United
\quellseite{190}
States. More than once John Wesley had written in praise of Raikes’ schools (for instance, under date of July 18, 1784: “I find these schools springing up wherever I go. Perhaps God may have a deeper end therein than men are aware of. Who knows but that some of these schools may become nurseries for Christians”). Every Methodist pastor, in the first \emph{Book of Discipline of Methodist Episcopal Churches} (1784), was directed “where there are ten children whose parents are in the society, to meet them at least one hour every week.” In 1786 Francis Asbury, the first bishop of the Methodist Church to be ordained in the United States, organized a Sunday School in the home of Thomas Crenshaw in Hanover county, Va. In 1790 the Methodist Conference at Charleston, S. C., ordered that these schools should be established “in or near the place of worship,” provided for the appointment of teachers, and set the sessions “from six in the morning till ten and from two o’clock in the afternoon till six, when it does not interfere with public worship.” The next year (1791) a Universalist Sunday School was established in Philadelphia, a Friends’ school in the same city in 1792, and a Baptist school at Pawtucket, R. I., in 1797. At an early time organizations were formed which had as their purpose the establishment of new Sunday Schools. On basis of Brown’s investigations, Cope furnishes the following summary: “The organization of the “First-Day or Sunday School Society” of Philadelphia, in 1791, was followed in 1808 by the organization of the “Evangelical Society” to promote Sunday evening schools in Philadelphia; the Pennsylvania Union in Pittsburgh, 1809; the Female Union Society in New York, 1816; the New York Sunday School Union, 1816; the “Boston Society for the Moral and Religious Instruction of the Poor,” 1816. All these were of local influence only, but a wider scope was designed for the “Philadelphia Sunday and Adult School Union,” 1817, for it planned to promote
\quellseite{191}
the organization of new schools in villages and the country. In 1821 it employed what was probably the first American Sunday School missionary; he organized sixty schools in six states. Doubtless this work prepared the way for the American Sunday School Union.” It is quite evident that the American Sunday School Union has adopted the chief principles of the Philadelphia Union and continues to apply them to this day. One of the most important of these was the idea that the Sunday School should be interdenominational. Thus we read in the Statutes of the Philadelphia Sunday and Adult School Union that one of its objects was “to cultivate unity and Christian charity among those of different names,” and in presenting their appeal to the public for generous support its leaders said: “The comparative fewness of Christians calls for all practicable and profitable union among themselves. Divide and conquer is the maxim of their great foe. Unite and triumph be then the motto of Christians.” Tracts in large numbers were issued to promote the formation of Sunday Schools and S. S. Associations, for instance a “\emph{Model Constitution},” a “\emph{System for the Internal Regulation of Sunday Schools}” including the “grading of the school in four grades, providing for two sessions of the school each Sunday, and thorough rules with a complete scheme of rewards and penalties for pupils, teachers, and officers, based upon record of attendance, lessons, and behavior.” In 1827 the Union issued 25,000 Sunday School Hymn books, 8000 class books, 2000 \emph{Teacher’s Guides} by John Angell James, 81,000 premium books, and for the third time, half a million of red and blue tickets.

The ever expanding work of this Philadelphia Union led to the formation, in 1824, of the American Sunday School Union, an organization that was national in name and scope. The objects of this Society were stated by the founders thus: “To concentrate the efforts of Sabbath-school associations
\quellseite{192}
in different sections of the country; to strengthen the hands of the friends of pious instruction on the Lord’s Day; to disseminate useful information; to circulate moral and religious publications in every part of the land; and to endeavor to plant a Sunday School wherever there is a population.” The basis of this Union was declared to be “No sacrifice of principle essential to salvation; no compromise of duty; no interference with the internal management of smaller associations; all discordant elements must be banished, and (there must be) union with Christ, and union with each other.” Upon this basis, the essential truths of Christianity held in common by all evangelical denominations, were to be fully inculcated. The individual members of the A.S.S.U. could maintain the integrity of their denominational relations, but in this work they were to unite with members of other churches for the purpose of teaching “the truths that Christ taught and as plainly as he taught them.” The significance of the A.S.S.U. for the development of the Sunday School can hardly be overestimated. It published a vast amount of literature for pupils and teachers, prepared lesson plans and recommended them to all schools, issued Bibles, text books, and other lesson helps down to the most modest text cards, and promoted the organization of new schools with brilliant zeal and enthusiasm. At the meeting in 1830 it resolved “within two years to establish a Sunday School in every destitute place where it is practicable throughout the valley of the Mississippi”; at once 78 missionaries started out from Cincinnati and founded 2876 Sunday Schools in the Mississippi valley during the next two years. In 1832 the A.S.S.U. invited representatives of the S. S. cause from all parts of the country to attend a meeting in New York; thus the first National Convention was brought about (October 3). An exhaustive questionnaire which had been sent out furnished valuable data concerning the actual conditions and
\quellseite{193}
needs of the S.S. On the basis of this material, Cope summarizes the situation at that time in the following manner: “(1) There were several distinct types of schools; infant schools, mission schools, adult schools, evening schools, as well as the general schools. (2) Lesson plans were many. Some spent over an hour in class memorization of long passages of Scripture; others simply told Biblical and other stories; some assigned one verse for each day of the week, the whole being the basis of recitation and comment on Sunday; the tendency was to adopt the limited portions designed by the A.S.S.U., all schools aided by this society being required to use the same lessons; everywhere the emphasis was on Biblical material. (3) Nearly all schools were directed by superintendents. (4) Libraries of general literature were established in many schools.\footnote{No public libraries existed in the early part of the nineteenth century, nor was there any juvenile literature; President Humphrey of Amherst could not recall over six books in this class, even counting Webster’s \emph{Spelling Book}, \emph{Robinson Crusoe}, and \emph{Pilgrim’s Progress}! The first S.S. library was established at Boston in 1812; in 1830 the A.S.S.U. issued 200 books for S.S. children.} The type of book afterwards to be known as a “Sunday School” book was almost a special creation for them. (5) Special “Bible classes” for church members and “adult classes” were organized. (6) Many schools met on Sunday afternoon, the sessions often being two hours in length. (7) In the frontier regions schools were organized before churches and became the parents of the latter. (8) In the cities there were many mission schools not immediately connected with specific churches.” This first national convention also resolved, “that the S.S. should embrace all classes of the community” and not only destitute and neglected children. The second national convention met in Philadelphia in 1833. It urged the establishment of Sunday Schools in jails and reformatories and the promotion of groups for the study of the Bible in the homes.
\quellseite{194}
For several reasons no national convention was held during the ensuing twenty six years. One important reason was the reawakening of denominational consciousness in this period: the more vigorously this consciousness was aroused, the more inevitable was the conflict with an association which was avowedly un-denominational and virtually super-denominational, which professed to preach a gospel of Christ free from all denominational “stains and encumbrances.” and which arrogated to itself some of the Church’s inalienable functions. So it came about that the various denominations organized their own S.S. Societies. In 1827 the Methodist Episcopal Church created its S.S. Union, the Unitarians followed in the same year, the Lutherans (v.i.) in 1830, the Congregationalists in 1832, the Baptists in 1840 (resp. 1826), the Southern Baptists in 1857, the Presbyterians organized, in 1833, “The Board of Managers of the Presbyterian Tract and Sabbath School Book Society” which subsequently promoted the S.S. work of this denomination. The reawakening of confessional consciousness thus made these 26 years a period of battle against the super-denominational claims and tendencies of the A.S.S.U. The historians of the A.S.S.U. contemptuously designate this movement as “sectarian”; as a matter of fact, however, this opposition was not only justified, but clearly necessitated by the very nature of the Church. While the A.S.S.U. thus was seriously limited in many of its activities, it devoted itself with redoubled energy to the publication of periodicals, teachers’ manuals, and S.S. literature in general and also to the formation of Sunday Schools in frontier regions where no separate denominational schools had been organized. Of special importance was the publication of the \emph{Sunday School Journal} (since 1831) in which the pedagogic theories of such great educators as Sturm, Comenius, Milton, Francke, Rousseau, Pestalozzi, Froebel were set forth and the current systems of Lancaster, Bell, Gall,
\quellseite{195}
Stow, and Mimpriss were examined as to their applicability to the S.S. Naturally the method and the standards of teaching were greatly improved by these efforts; also the denominational schools profited thereby. Another tremendously important service rendered by the A.S.S.U. to the Sunday School and the cultivation of a Christian spirit in general was the creation of a juvenile literature of a religious type. According to E. W. Rice, the following principles governed this phase of the Union’s work: “(1) The literature must be suited to the development of the child mind; (2) pure in tone; (3) serious rather than sensational; (4) ingenious, but not absurd; (5) popular rather than polished; (6) thoroughly Biblical and evangelical; (7) truly American, written by American writers, statesmen, and philanthropists, redolent of American mountains, forests, prairies, rivers, history, and yet filled with the spirit of the Word, not of the world.” In 1845 the Union began to issue a one-hundred bound-volume (72 to 272 pages per volume) library for Sunday Schools at the low price of ten dollars. Three similar sets followed. Add to these publications the periodicals for children (\emph{The Youth’s Penny Gazette}, with full page illustrations, was issued from 1843), the hymnals for which men like Mason and Hastings composed melodies, the large number of new Sunday Schools founded by its representatives, and the arousing of missionary zeal (as early as 1835 the Union declared one of its objects was “to assist in carrying the Gospel to every family in the world, and to insure the religious instruction of every child that is born”)—and one readily sees what a rich stream of blessing proceeded from the A.S.S.U. even after its activity had been circumscribed by the newly founded denominational schools and societies. Nor should it be forgotten that many denominational societies adopted the lists of lessons issued by A.S.S.U. and explained them from their own point of view.
\quellseite{196}
In 1859 the third national convention was in session at Philadelphia. On this occasion also the western section of our country was well represented; the most prominent men besides Henry Clay Trumbull (Congregationalist, since 1871 secretary of the A.S.S.U. and chairman of the executive committee of the National S.S. Convention, soon afterward editor and part-owner of the \emph{Sunday School Times}, died 1903) were Benjamin Franklin Jacobs of Chicago (a Baptist business man, prominent S.S. leader in Illinois, conceived the idea of Uniform Lessons, since 1872 active member of the Lesson Committee, later president of the International S.S. Association, died 1902) and John Heyl Vincent of Joliet, Illinois (Methodist, one of the leaders in the movement for the International Uniform Lessons, helped found the Chautauqua Assembly (1874) where training courses for S.S. teachers were offered, Bishop since 1888). After the Civil War the fourth national convention met at Newark, N. J., in 1869; and at the fifth convention at Indianapolis, in 1872, the plan of Uniform Lessons was adopted, a Lesson Committee (composed of twelve members, later increased to fourteen by the addition of two Canadians) was appointed, and the national convention was made International with biennial meetings. In 1889 the first World Sunday School Convention was held in London, chiefly owing to Jacob’s efforts. At the fifth convention in Rome (1907) it changed its name to World’s Sunday School Association. The purpose, policy, and field of the Association were defined as follows: “That this Association shall hold conventions and gather information concerning the conditions of Sunday Schools throughout the world by correspondence, visitation, and other methods; that it shall seek to extend the work and increase the efficiency of Sunday Schools by cooperation with Sunday Schools and missionary organizations and otherwise, especially in those regions of the world most in need of help; that it shall seek
\quellseite{197}
to improve, so far as possible, the methods of organization and instruction in Sunday Schools and promote the formation of S.\,S. Unions and Associations.” By agreement the world field was divided for the purpose of financial administrative responsibility: Europe, Australia, South Africa, India—the British Section of the Executive Committee; North and South America, Japan, Corea, Philippines, Turkey, North Africa—the American Section of the Committee; China to be jointly administered by the British and American Sections. The chief feature of the seventh convention at Zuerich, Switzerland, in 1913 was the report of the Six Commissions of the Association on the Sunday School conditions, needs, and opportunities of the world, covering continental Europe, South Africa, India, the Orient, Latin America, and the Mohammedan lands. According to the latest available statistics there were in the United States 143,035 Protestant Sunday Schools with 17,729,216 pupils. To the World Association belonged in 1924 32,677,611 members in 347,001 schools with 3,520,192 teachers and officers.

Efforts are made to introduce everywhere within the International S.\,S. Association the identical lessons which are interpreted independently by denominational boards according to their convictions. In this way the autonomy of the cooperating denominations is preserved and at the same time expression is given to the unity of the Churches composing the Association—a compromise terminating the conflict between the superdenominational aspirations of the A.\,S.\,S.\,U. and the denominational demands. Most of the Churches have adopted these lessons, but there are many Sunday Schools numbered in the statistics of the I.\,S.\,S.\,A. (e.\,g. most of the Lutheran schools) that use their own systems. In order to do justice to the common tasks of all Sunday Schools (especially the improvement of methods) more fully than is possible at the International Conventions, “\emph{The Sunday}
\quellseite{198}
\emph{School Council of Evangelical Denominations}” was organized in 1910 which included in its membership about 150 S.\,S. editors, publishers, and secretaries representing all the more important evangelical denominations in the United States and Canada. The chief object of this Council was to apply the results of modern psychology, pedagogics, and sociology to the S.\,S. lesson courses without interfering with denominational autonomy in matters of administration. Its resolutions though helpful and salutary now and then, on the whole evinced unionistic tendencies shifting the emphasis from the center of the New Testament. This Council and the I.\,S.\,S.\,A. were merged into a new body in February 1922, “\emph{The International Sunday School Council of Religious Education}.”

In reference to the study material used by American Sunday Schools it may be said that in general the Bible formed the basis of instruction ever since the S.\,S. did not need to devote its time to the study of elementary branches. Only a few recent lesson plans depart from this practice, offering general materials (natural history, tales, daily life, etc.) which are to be interpreted and applied religiously. In regard to the method employed in introducing the pupils to the Bible, several periods stand out more or less distinctly. At first the memory method was deemed sufficient. The pupils memorized and recited shorter or longer sections of Scripture. Spurred on by rewards, some children would learn a hundred or more verses, or even a whole Biblical book in the course of one week. Now, to store in the mind a rich treasure of Bible passages is undoubtedly beneficial; but when the entire S.\,S. period, even when it lasted two hours, was consumed by the recitation of memory material, no time was left for explanation and application, and very little religious insight was brought about and still less influence was exerted upon the children’s religious life. A step of progress was taken
\quellseite{199}
when in 1825 Gall’s method was introduced in the Sunday Schools of America. Dr. James Gall of Edinburgh, Scotland, had prepared a series of Bible lessons on a method called “\emph{Nature’s Normal School}”; the lesson consisted of a brief Bible story to be told by the teacher and a series of questions on that story. The system, however, was so complicated and mechanical that little actual knowledge was imparted and the soul of the pupil remained cold. In 1829, Rev. Albert Judson projected for the A.\,S.\,S.\,U. a scheme of lessons to be completed in five years; all classes studied the same lesson and reviews were conducted at the end of the quarter. This plan was adopted by the A.\,S.\,S.\,U. and widely used during the following years. In distinction from the Uniform Lessons introduced in 1872, the period inaugurated by the acceptance of Judson’s Lessons may be designated as the \emph{First Uniform Lesson Period}. This method actually provided for the study of the Bible in the course of five years, called for the actual reading of the Bible, and stimulated the children to think about the stories for themselves. But, of course, everything of a denominational nature was omitted. Among the numerous helps for these lessons, the most popular were \emph{Judson’s Questions} in three grades, and especially the \emph{Union Questions} edited by John Hall (in five little volumes, later after an extension of the course in seven, and finally the review year having been added, in twelve). For beginners \emph{The Child’s Scripture Question Book} was issued, and for adult Bible classes McDowell’s and Tyng’s lessons. In point of method Cope’s judgment is: “The examination of a text book of 1830, with the assigned lesson, analysis of the narrative, series of questions, explanations of unfamiliar words, exposition on the doctrine and applications on practical lessons, shows that these ‘lessons’ were at least as good as those in the average ‘quarterly’ of sixty years later.” A great diversity of lesson plans and methods resulted from the organization
\quellseite{200}
of denominational S.\,S. associations for which reason the period following upon the first uniform lesson period is often designated as the \emph{Period of Babel Systems}. Many schools continued to use the Uniform Lessons; others devised their own; some sought to adapt the lessons to the age of the pupils (e.\,g., the Unitarians had a series of eight graded texts in 1852); Episcopal schools prepared lessons on the program of the church year. John H. Vincent published in his \emph{Sunday School Teacher} a course entitled “Two Years with Jesus.” The capable albeit somewhat passionate Edward Eggleston, in his \emph{National Sunday School Teacher} published another series of lessons in 1867 and in his \emph{Manual: a Practical Guide to Sunday School Work} showed with great skill how to use the plan. He was an outspoken foe of the Uniform Lesson (the same lesson for all classes on one Sunday) and was convinced that “the needs of schools call for a variety of lessons just as a day school needs more than one reader in Arithmetic.” At the national convention of 1869, however, the lack of uniformity was keenly felt and it seemed imperative to do something toward introducing uniform standards and methods in the S.\,S. At this juncture B.\,F. Jacobs of Chicago proposed his plan of Uniform Lessons for all Sunday Schools, the helps to be prepared by the denominational boards. Because the autonomy of the denominations was sufficiently safeguarded and at the same time the ideal of uniformity was, in a measure, approached, the Convention of 1872 with an overwhelming majority adopted this plan. Eggleston’s heroic opposition did not stem the tide (he declared that it was “a movement backward; it would pull down good schools, produce a dead level uniformity”; and in 1880 he said: “The system of rigid adherence to one lesson for all the school, combined with selections now and then of subjects fit only for a theological seminary, is not in accordance with practical wisdom”). An international lesson committee was
\quellseite{201}
appointed (see p. 174). Therewith the period of the \emph{International Uniform Lessons} began. For thirteen years the new lessons dominated the entire field; all opposition was quelled. When the preparation of a separate series of lessons for younger classes was demanded in 1885, the Lesson Committee refused to issue such lessons. The London S.\,S. Union in 1890 requested “a system of graded lessons, in which the same subjects shall be considered, but with different Scripture chosen first for the primary classes, and then for more advanced classes,” but the request received as little attention. Rev. Erastus Blakeslee even began to issue, at his own risk, a graded series of inductive Bible lessons which formed a six year course. The leaders of the I.\,S.\,S.\,A. were so enamored with the idea of uniform lessons that they did not react favorably to these suggestions and criticisms although it should have been evident that beginners and adults have special needs that can not be supplied by a mere gradation of the lessons helps, but only by different lesson texts and subjects. It is quite significant that in the following years the Primary teachers were the most persistent objectors to the Uniform Lessons and the most aggressive champions of a graded system. With increasing insistence they petitioned the Lesson Committee to issue suitable lessons for little children, especially since the interest in child-psychology began to grow in our country (at a conference in 1903 one of the speakers bluntly said: “The old education put the material first, the child second; the new education puts the child before the material!”). After many heated debates a conference of I.\,S.\,S.\,A. leaders finally adopted the following resolutions: “(1) That the system of a general lesson for the whole school which has been in successful use for 35 years, is still the most practicable and effective system for the great majority of the Sunday Schools of North America. Because of its past accomplishments, its present usefulness,
\quellseite{202}
and its future possibilities, we recommend its continuance and its fullest development; (2) that the need for a graded system of lessons is expressed by so many Sunday Schools and workers that it should be adequately met by the I.\,S.\,S.\,A., and that the Lesson Committee should be instructed by the next International Convention, to be held in Louisville, Ky., June 18–23, 1908, to continue the preparation of a thoroughly graded course covering the entire range of the Sunday School.” The convention at Louisville unanimously adopted this report. But hardly had the new lessons appeared when the Southern Baptists and the Southern Presbyterians justly attacked them for the “absence of doctrine, the presence of extra-Biblical lessons, the omission of many important topics and liberal interpretation of Scriptures.” These and other criticisms resulted in some slight improvements. Graded lessons have now been issued as follows: Beginners (four and five years of age) two years; Primary (six to eight years) three; Junior (nine to twelve years) four years; Intermediate (thirteen to sixteen years); Senior (seventeen to twenty-one years) four years. From 1908, accordingly, dates the period when alongside of the International Uniform Lessons the \emph{International Graded Lessons} exist by right in the I.\,S.\,S.\,A. Of the schools that have introduced graded lessons and consistently applied modern psychology and pedagogy (often, indeed, at the expense of Bible truth) and have exerted strong influence, for good and evil, we mention the Model Sunday School of Columbia University, New York City; University Congregational School of Hyde Park, and the Hyde Park Baptist Sunday School, both of Chicago and in close touch with the University of Chicago.

\trenner

The Lutheran congregations in America were in the early period completely dependent on the catechetical work of the home Church. In his plea for help of May 31, 1693,
\quellseite{203}
Springer also asked the home Church in Sweden for 300 catechisms. The first German congregations were closely connected with the pietistic circles in the homeland. As a case in point, Boltzius of Charleston reports in 1736 how he explained the catechism in the daily evening service; we read the following entry: “Last evening we reviewed the seventh commandment and catechized upon the eighth. The consciences of some in the congregation were deeply stirred. In the exposition of the commandments to simple hearers we find it highly necessary to indicate the sins forbidden and the virtues enjoined therein not only generally, but they must be clearly specified according to the circumstances of the hearers.” Muehlenberg’s fidelity in the instruction of the neglected youth is well known. With more regularity than was possible for Muehlenberg who had to travel about extensively, Brunnholtz whose infirmity kept him closely to his parish, conducted \emph{Kinderlehre} or “Children’s Instruction.” The younger children he instructed in the Small Catechism, and the older ones and the servants in the “\emph{Order of Salvation}” and Bible History, an arrangement which clearly indicates Francke’s influence. It is to be regretted that, later, Rationalism entering from Germany and other countries made inroads upon American church life and affected also the catechetical text books. Quite undisguised it appears in F.\,G. Quitman’s \emph{Evangelical Catechism} (Hudson, 1812, or possibly 1814). Opposition to this perversion was at times rather feeble, but it never subsided altogether and kept alive the Lutheran confession in catechetical instruction.

In 1787 Professor John Caspar Velthusen of Helmstaedt prepared and edited the \emph{Lehrbuecher fuer die Jugend in Nord Carolina} (Text-books for the youth of N.\,C.) at the request of Nuessmann who had been sent to care for the Lutherans in Mecklenburg county in North Carolina. The purpose was to publish (1) a catechism, (2) a book of questions on the catechism, (3) a Bible handbook for everyone, (4) a selection of Bible stories with a brief history of religion, (5) a
\quellseite{204}
book of useful information for every day, (6) a manual of civics, (7) a manual of geography. At least the first two of these projected books appeared and are extant, viz., \emph{Helmstaedtischer Katechismus oder Christlicher Religionsunterricht nach Anleitung der heiligen Schrift}, and \emph{Fragebuch fuer Eltern und Lehrer, oder Anleitung zu Fragen und Gespraechen ueber den Katechismus mit Ruecksicht auf die Verschiedenheit der Faehigkeiten und des Alters der Jugend}. In his catechism which is a worth-while selection of Bible passages, Velthusen adheres to the doctrine of the Trinity, the eternal Sonship of Christ, the virgin birth, and the vicarious atonement, but in his statements concerning the appropriation of salvation he pays tribute, in terminology and subject matter, to the prevailing rationalism.—To this period belong also the following catechetical treatises written from widely diverging viewpoints: H.\,E. Muehlenberg, \emph{A Companion to the Catechism, or a Course of Instruction in the Christian Religion, for the Benefit of the Young}. A. Wackerhagen, \emph{Inbegriff der Glaubens- und Sittenlehre}, Philadelphia, 1804. Anon., \emph{Anfangsgruende der Religion, oder Katechismus fuer kleine Kinder}, Lebanon, 1814. G. Vorberg, \emph{Handbuch fuer die Konfirmierten}, New York; in English translation with the title \emph{Way Marks for the Confirmed}, 1815. Ph.\,F. Mayer, \emph{Instructions in the Principles and Duties of the Christian Religion, for Children and Youth}, 1816. Paul Henkel, Deutscher und englischer Katechismus. J.\,G. Lochmann, Hauptinhalt der christlichen Lehre. Id., Evangelical Catechism, Harrisburg, 1882. E.\,L. Hazelius, Materials for Catechization on Passages of Scripture, 1823.

The educational institutions available in the Lutheran Church were the parish and catechumen schools (the Pennsylvania Synod alone had as late as 1812 one hundred and sixty parish schools) and, in many congregations, instruction of the young on Sundays (Christenlehre). When the Sunday School, in its new form, was introduced in America, the Lutherans, too, soon adopted it as an additional agency of instruction. German Lutheran congregations were the pioneers in this respect. On May 7, 1804 a widow organized a S.\,S. of five or six children in St. Michael’s and Zion’s Congregations in Philadelphia; St. Paul’s Church followed in 1805. The North Carolina Synod in 1811 and again in 1812 called on the congregations to form Sunday Schools.
\quellseite{205}
In 1819 Christ Church, York, Pa., and Zion’s Church, Harrisburg, Pa., organized Sunday Schools. Adam Keller’s efforts resulted in the formation of a S.\,S. in St. John’s Church at Philadelphia in 1821 which on the opening day had an attendance of 19 teachers and 60 children in the forenoon and 100 children in the afternoon. Schools came into being at Lancaster in 1828 and at Reading in 1829. Progress was slow chiefly for the reason that the considerable number of day schools still in existence lessened the need for the Sunday School in its new form. The Pennsylvania Synod had 26 schools in 1838, in 1844 already 73. Six years after the formation of the American Sunday School Union, the General Synod organized a Lutheran Sunday School Union (1830). Influential men sponsored the S.\,S. cause, e.\,g., S.\,S. Schmucker (\emph{Plea for the Sabbath School System}, delivered at the anniversary of the Gettysburg S.\,S., 1830), S.\,A. Seiss (\emph{Thoughts on Education}, address to the Sabbath School Union, Cumberland, Md., 1850), and Benjamin Kurtz (\emph{Sermon on Sabbath Schools}). When it became evident that the parish school would not thrive in the Church as far as she used the English language, and when the inadequacy of pre-confirmation catechization was recognized, the S.\,S. became a practical necessity. “Father” F.\,C. Heyer, best remembered as America’s pioneer Lutheran missionary to India, who in his youth had been parish school teacher in Southwark, traveled very extensively as the agent of the Lutheran S.\,S. Union and established a large number of Sunday Schools in the thirties of the last century. More rapid progress was made since about 1840, the year in which such large congregations as St. Matthew’s and St. James’ in New York established schools. In German quarters it was especially Brobst who promoted this work in view of the horrifying ignorance of many children even in church homes. He issued a booklet for the instruction of beginners, and also showed
\quellseite{206}
how to open and to close the Sunday School session liturgically. In 1854 he was appointed, together with Vogelbach and B.\,M. Schmucker, to address a circular letter to the congregations of his Synod in the interest of the S.\,S. He accomplished his task so perfectly that an English translation of this letter was published as late as the eighties. Lutheran hymnals and lesson leaves introducing into Scripture and evangelical doctrine were issued, and the Catechism was adapted to S.\,S. use. As long as the instruction of confirmands was faithfully performed, Luther’s Small Catechism was used which had been translated into English long ago (of special merit are the translations made by Ph.\,F. Mayer, 1816; by C.\,F. Schaeffer, made in 1895 by order of the Pennsylvania Synod; and the Joint Committee’s translation, 1899); also new catechisms were produced. Of the expositions of Luther’s Catechism, Mann and Krotel’s work stands out as especially valuable (issued 1863 by order of the Ministerium of Pennsylvania, German and English, based on Caspari’s exposition). The \emph{Augsburg Sunday School Teacher} served as a suitable teacher’s manual; and periodicals were issued for children and for Sunday Schools. H.\,J. Schmidt (then professor in Gettysburg College and Seminary) published a \emph{History of Education} in 1842 which contained in part “a plan of culture and instruction based on Christian Principles, and designated to aid in the right education of youth, physically, intellectually, and morally”; it was the only work of the kind which the American Lutheran Church had produced till then. And it remained the only work until H. Ziegler’s \emph{Catechetics, Historical, Theoretical, and Practical} appeared in 1873. Sunday School teachers’ meetings are first mentioned in 1852; B.\,M. Schmucker conducted such meetings weekly for the purpose of studying the lesson for the following Sunday.

Of the catechetical works that appeared during this period we mention the following: John G. Morris, \emph{Catechumen’s and Communicant’s}
\quellseite{207}
\emph{Companion}, Baltimore, 1831. Id., \emph{Catechetical Exercises on Luther’s Catechism}, adapted from the German, Baltimore, 1832. S.\,W. Harkey, \emph{Lutheran Sunday School Question Book, or a Help to the Systematic Study of Sacred Scriptures}, compiled from the German, Frederick, Md., 1838. John G. Morris, \emph{Popular Exposition of the Gospels for Families, Bible Classes, and Sunday Schools}, 2 vols., Baltimore, 1840. Benjamin Kurtz, \emph{The Serial Catechism, or Progressive Instruction for Children, adapted to the growth in grace and knowledge}, Baltimore, n.\,d. Id., \emph{Easy Catechism for Young Children}, 1843. John G. Morris, \emph{Luther’s Catechism Illustrated by Additional Questions and Answers}, 1844. C.\,A. Morris, \emph{An Easy Catechism for Young Children}, n.\,d. C. Phil. Krauth, \emph{Lutheran Sunday School Hymn Book}, about 1855. S.\,A. Seiss, \emph{Training Little Ones for Christ}, Baltimore, 1853. Clemens Miller, \emph{Luther’s Catechism}, Baltimore, 1857. M. Sheeleigh, \emph{Hymns, Selected and Original, for Sunday Schools of the Evangelical Lutheran Church, Revised, with Appendix, containing Hymns for the Use of Infant Schools}, Baltimore, 1860. T.\,T. Titus, \emph{Sunday School Question Book}, n.\,d. S.\,S. Schmucker, \emph{Evangelical Lutheran Catechism, Designed for Catechumens and the Higher Classes in Sabbath Schools}, Baltimore, 1850. W.\,I. Mann and G.\,F. Krotel, \emph{Luther’s Small Catechism, Explained in Questions and Answers for the Use of the Church, School and Family}, 1863. M. Sheeleigh, \emph{Outlines of Old Testament History}, also \emph{Outlines of New Testament History}, about 1865. E.\,I. Koons, \emph{Questions on Luther’s Catechism}, 1869. G.\,F. Jager, \emph{Katechismus der christlichen Lehre in Fragen und Antworten}, Kutztown, 1833. H.\,W. Scriba, \emph{Anfangsgruende des Christentums}, translated from the French, Chambersburg, 1834. J.\,G. Schmauk, \emph{Erstes Buch fuer deutsche Schulen}, Philadelphia, 1844. F.\,P. Peixoto, \emph{Leitfaden, wonach der kleine Katechismus Luthers erklaert wird}, Sunnytown, Pa., 1845. S.\,K. Brobst, \emph{Fragebuechlein fuer die Anfangsgruende der Religion}, 1846. Id., \emph{Gebete fuer Sonntagsschulen}, 1847. Id., \emph{Fragen und Antworten aus dem Neuen Testament}, 1849. C.\,A. Morris, \emph{Liederbuechlein fuer Kinder}, gesammelt und herausgegeben von den Lehrern der deutschen Sonntagsschulen in Yorktown, 1849. F.\,C. Wyneken, \emph{Spruchbuch zum kleinen Katechismus Lutheri}, Baltimore, 1849. S.\,K. Brobst, \emph{Gesangbuch fuer die Sonntagsschulen der ev.-luth. und deutsch reformierten Kirchen}, Allentown, Pa., 1853. J.\,C. Haas, \emph{Luthers kleiner Katechismus mit Spruechen aus der Heiligen Schrift}, Philadelphia, 1856. Mann u. Krotel, \emph{Luthers Katechismen mit Fragen und Antworten}, 1863. C.\,F. Spring, \emph{Katechetische Unterweisung zur Seligkeit}, New York, 1868.

\quellseite{208}
When the lessons issued by the International Sunday School Association began to appear in 1872, they were introduced more or less universally in the English Lutheran synods of the East, especially in the General Synod where they were explained in the Augsburg Teacher and in the church papers. Not even the General Council was altogether proof against them. Quite generally catechumen instruction was retained as a supplement for the Sunday School. The text books chiefly used were Mann and Krotel, \emph{Luther’s Small Catechism}; also Seiss, Spaeth, and Jacobs, \emph{Luther’s Small Catechism}, with Scripture Texts, by authority of the General Council in America. Philadelphia, 1883; or the Stohlmann reprint of Michael Walther’s Catechism as remodeled by Luehrs (the English translation by H.\,E. Jacobs as well as the German edition were widely used). Other serviceable books were Wischan and Spaeth, \emph{Mein erstes, zweites, drittes Sonntagsschulbuch}; the Biblical History published by the Pilger Book Store of Reading; F. Greenwald, \emph{Questions on the Gospels and Epistles of the Church Year}, Lancaster, 1873. 1874; Ludwig’s abovementioned catechism with proof texts and an appendix containing The Order of Salvation; \emph{Sunday School Book of the General Council} (English 1873, German 1897, revised 1906); Wonneberger, \emph{Sonntagsschulharfe}.

The inadequacy and the imperfections of Sunday School instruction as based upon the International Uniform Lessons were keenly felt in the General Council. In the first place, the I.\,S.\,S.\,A. encourages the view that the S.\,S. is a separate, independent institution alongside the Church; furthermore, the fact that the children are baptized is not properly recognized; little, if any, real knowledge of the history of salvation is imparted; law and gospel are so often confounded; the central facts of salvation are not sufficiently clearly brought out; no instruction in the catechism being provided for, the real goal of instruction is obscured; the appreciation
\quellseite{209}
of specific catechumen instruction is suppressed rather than aroused; and finally, an organic connection of the instruction of the confirmed with the lessons learned during their pre-confirmation instruction is rendered impossible. Moreover, the Uniform Lesson System does not take sufficient account of some fundamentals of psychology and pedagogics. It was worthy of the highest praise, therefore, when the General Council in the face of severe opposition decided in 1895 to supply the urgent need of a typically Lutheran and pedagogically sound lesson system by publishing its own system. A more capable and purposeful editor than Dr. Theodore Schmauk could hardly have been found. It is true that a graded system had been advocated before (cf. the preceding sketch of the development of the S.\,S.), and some preliminary work had been done; but this was the first time that a Lutheran S.\,S. system was created and that a whole church body united in this sort of enterprise. The system is divided into Primary, Intermediate, and Senior Departments. The Primary Department is divided into Kindergarten, Children’s, and Junior Department (from the third to the ninth year); \emph{Wonderland and Sunbeams}, \emph{Workland and Sunshine}, \emph{Pictureland and Sunrays} are the lesson helps for this department. \emph{In Mother’s Arms} and \emph{At Mother’s Knees} are preliminary to the whole, serving the mothers in instructing their little ones at home. The Intermediate Department is composed of seven grades (from the tenth to the sixteenth year); teaching aids to be used are \emph{Bible Story, Bible Readings, Bible History, Bible Facts and Scenes, Bible Biography, Bible Teachings, Bible Literature}. Instruction in the Catechism with confirmation is to occur between Bible Teachings and Bible Literature, or rather run parallel with these grades. The Young People’s Department (17 to 20 years of age) and the Adult Department (21 years and above) devote themselves, with the aid of the \emph{Senior Lesson Book}
\quellseite{210}
or Commentary, to the study of Scripture proper. The General Synod is still connected with the International S.\,S. Association, using its Uniform, or Graded Lessons; but in 1911 this body authorized its S.\,S. Committee to enter into negotiations with other English speaking bodies for the purpose of creating in cooperation with them a specific Lutheran Sunday School literature. In the East there is used for the instruction of confirmands, in addition to the literature already mentioned, the \emph{Cathechism} by the author of this text book on Catechetics; Loehe’s catechism, translated by Horn; and the catechisms by Trabert and Stump. On the basis of Kaftan’s work (cf. § 19), I.\,W. Horine, in his \emph{Catechist’s Handbook}, Philadelphia, 1909, has given the catechist a valuable aid to the understanding of Luther’s Catechism.

In the third decade of the last century German, and a little later also Scandinavian, immigration into the western part of our country commenced on a large scale. The newly founded congregations and church bodies had brought some agencies of religious instruction with them from their mother countries, others they had to devise for themselves. Catechumen instruction and Christenlehre they brought with them; among the new agencies, devised in this country, is the parish school. In the public state-supported school which every child in the home country must attend, religious instruction, indeed, occupied a prominent place, but the school was a state school, not a parochial school; and if the Church exerted any influence upon religious instruction as imparted in the public school, it was only by reason of being a State Church. In their new home, however, the immigrants found Church and State—fortunately—separated by constitutional enactment; the State, therefore, could not possibly impart religious instruction in school. But for the sake of her self-preservation the Church could not dispense with religious instruction of her youth; and since education ought to be a harmonious
\quellseite{211}
unit, the spirit of Christ penetrating the sum of secular knowledge, there arose for these new congregations and church bodies the necessity of establishing independent church, or parish schools. In these schools they gathered their children and endeavored to train them by graded instruction in religious and secular branches to become efficient members of the Church and of the State. It was principally Loehe who advocated the formation of such schools; he also collected funds for the first Lutheran Normal School (Saginaw City, Mich., 1852) and sent the teachers and students for this school. The Synod of Iowa founded by Loehe’s disciples has never quite forgotten this task; but the course of developments caused the Missouri Synod to become the chief exponent and champion of the parish school ideal. Missouri’s care for the parish school is and remains the most splendid chapter in its history. In Lindemann’s \emph{Schul praxis} its principles of school management were laid down; its pedagogical organ, \emph{Schulblatt}, edited by Lindemann, has outlived all similar periodicals; and in the preparation of schoolbooks and other helps it has usually been guided by sound pedagogical principles. It is gratifying that the Missouri Synod at an early stage of its development not only realized the necessity of cultivating the language of the country in its schools, but even of founding purely English parish schools. In 1926 Missouri had in its two Normal Schools (River Forest, Ill., and Seward, Neb.) 630 students. Throughout the Missouri Synod there are 1388 parish schools with 1272 male teachers and 447 female teachers and 401 pastors teaching school. The pupils number 80,173. There remain to be mentioned the Normal School at Waverly (Iowa Synod), and at New Ulm (Wisconsin and Minnesota Synods). The Scandinavians have normal courses in connection with several of their institutions. Where conditions are adverse to the formation of complete parish schools, faithful congregations organize “partial” parish schools (cf. § 31); for instance,
\quellseite{212}
the pastor conducts a winter or confirmation school; there is frequently a summer school term of two months and Saturday School; and almost universally we find the Sunday School. The last named institution has made its way also in the West: at present every synodical body has a special Sunday School Committee which provides for the needed Sunday School literature. After years of purposeless efforts sound principles are gradually beginning to find recognition. The Sunday School systems of the Norwegians and of the Iowa and the Missouri Synods deserve specific mention.

It remains to note, briefly, the catechetical literature devised by the western Synods. Stephan and Walther introduced from Saxony the \emph{Dresden Cross Catechism} composed in 1688 by the ministerium of the Cross Church in Dresden. This was later generally replaced in the Missouri Synod by the large and small catechism of C. Dietrich (first edition 1613). The disciples of Loehe, in the Iowa Synod particularly, disseminated his catechism (first edition 1845) which had been composed with a special view to the Lutherans in America. Also the Swedes, Danes, and Norwegians made use of the catechisms of their respective homelands, e.\,g., of that of Pontoppidan. At present almost every Synod has its own synodical catechism although some trans-atlantic catechisms are still in use. Almost invariably they are bilingual. The catechism of the Missouri Synod has been edited by Schwan; that of the Wisconsin Synod was based upon the \emph{Dresden Cross Catechism}, but has been replaced by one edited by Gausewitz; that of Ohio is based upon Dietrich’s, but is about to be revised. Since 1904 also the Iowa Synod has its own catechism, edited by the author of this textbook on Catechetics. It differs from other catechisms, as to contents, in that it returns to the expository principle of Loehe (1845) and Tetelbach (1568); it excludes all supplementary material taken from dogmatics or sacred history, and it restricts
\quellseite{213}
itself to those subjects which are essential to a Christian life; as to form, in that it adopts the thetic, not erotematic, presentation. Since 1917 it is to be found, in both German and English, also in abbreviated size and in the form of questions and answers. The Augustana Synod, likewise, has its own explanation; the Norwegian Synod uses largely Dietrich’s Catechism while the United Norwegian Church employs the explanation by Sverdrup.

Among the textbooks on Biblical History those of Reading (Pilger Book Store, Reading, Penn.), the old Huebner Bible History (1714) and those issued by the Missouri Synod enjoy a large circulation. Among those in English, that published by the Augustana Synod deserves special mention; the latest, with several novel features, is by the author of this textbook. The book \emph{How I Tell the Bible Stories to My Children}, (second, revised edition, two volumes, 1926) by the same author, is designed to furnish a contribution to the art of story-telling as applied to Biblical History. Among aids for the catechist there are helps for the instruction in Biblical History by Wegener and by Simon in German, and by Rupprecht in English; for instruction in the Catechism by Stellhorn in German and Schuh in English. We mention also the writings of Schaller and Reu on instruction in the Bible (both English). In the Missouri Synod Quarterly Professor Dau has published material relative to instruction in the Catechism. Among theoretical dissertations there are to be mentioned, in addition to Lindemann’s Schul praxis, especially Grossmann’s \emph{Gemeindeschule}, Herzer’s \emph{Katechetik}, and Gerberding, \emph{The Lutheran Catechist}; also the book of the Methodist Hertzler contains some good sections. A sourcebook by the author in ten volumes is evidence that the Lutheran Church in America has taken part in historic research work in the field of Catechetics.

\trenner

\quellseite{214}
The efforts to introduce \emph{moral and religious instruction into the public schools} can not and shall not receive notice here. The publications of the Religious Education Association (headquarters in Chicago) contain considerable material in regard to this subject, as in regard to the religious training of our people in general. Some of the efforts, however, to complement public school education by religious day schools ought to be mentioned. What Charles De Garmo says in \emph{Principles of Religious Education} (p. 63): “From the standpoint of the development of religious intelligence the American (public school) system must be pronounced the most fragmentary, partial, inefficient, haphazard system in the world,” is more and more recognized as true; likewise, that the work of the Sunday School cannot be recognized as an adequate complement. G. Boville, of New York City, has briefly put the matter as follows: “The Church has enrolled about 15,000,000 children of school age in Sunday Schools, for whom are provided fifty-two half hours of instruction; there are 10,000,000 more children of school age not enrolled in any Sunday School.” In order to give the Sunday School pupils more religious instruction, and the others at least a little, this man gathered during the summer vacation the children of an East side precinct in New York City every day for six weeks in such churches as were opened to him and engaged college students to supervise and instruct them, and thus “brought together idle children of the streets, idle buildings of the churches, and idle students of the colleges.” Instruction, in part at least, was of a religious character “on broad non-sectarian lines.” In view of the success of Boville, the National \emph{Daily Vacation Bible School} Association was organized in order to arouse interest in this matter also at other places. In recent years the movement has grown tremendously. The association was made international and recently merged with the International Council of Religious
\quellseite{215}
Education. It is not saying too much that “a Vacation Bible School is now regarded as an essential department of every modern church program.” Already before this, another method had been tried in a small way. As early as 1898, H.\,R. Vaughn, pastor of a Congregational Church in Elk Mound, Wis., organized vacation religious day schools (two weeks) in conjunction with teacher training institutes which he conducted in the country and the smaller towns. He graded his schools on the basis of the public school and demanded credentials similar to those asked for by the public school boards, and with the help of these teachers sought to give the children instruction in Biblical History and in the history of missions and of the Church, and thereby to arouse in them enthusiasm for everything noble and good. Such schools — no propaganda for their extension was made — are conducted in a dozen villages and towns of Wisconsin, but also in some larger cities, e.\,g., Eau Claire, Madison, Beaver Dam, and Rockford, Ill. Professor W.\,J. Mutch of Ripon College is still the soul of the movement. In the summer of 1918 the Lutheran churches of Madison, Wis., conducted a Bible school which in some respects conformed to the movement fathered by Vaughn, but was governed by Lutheran principles. A kindred plan is prosecuted in Colorado, the so-called Greeley Plan of Bible Study. By virtue of cooperation between the State Teachers College of Colorado at Greeley and the churches, courses in religion given in the churches and approved by the College are accepted for credit. In 1907 the Lutheran pastor G.\,U. Wenner, as chairman of a committee appointed by the Federal Council of Churches of Christ in America, proposed another plan. He desired that the legislature “should authorize all children in the public schools to attend, at the option of their parents, on Wednesday afternoons, either the public schools or schools of religion in their own churches.” The salient point here
\quellseite{216}
is confessional instruction; and undoubtedly this plan marks a considerable advance, but the amount of instruction which it provides is still inadequate. — Confessional religious instruction during the week is made possible also by the \emph{Gary Plan} (as originally conceived) which W.\,E. Wirt, superintendent of the public schools of Gary, Ind., began to carry out in 1913. His desire was to care well for the child and not merely to give him formal instruction. He wanted more of the child’s time, so that the child might live most of his waking hours under the supervision of the school. “On this theory the Gary schools were conducted from 8:15 A.M. to 4:30 P.M. for six days in the week, and the school undertook to provide for the child’s play as well as his welfare in other lines. During these additional hours, Mr. Wirt thought the children should be sent in small parties to receive the benefit of any welfare agency, and among these stood the Church. Mr. Wirt informed the various churches in Gary that he stood ready to send to them for week-day instruction any and every child whose parent should sign a card desiring to have him so sent, and that the school would allot the church from one up to six hours per week of each child’s time, according as the church felt capable of using such time. This hour included the going from and returning to school. There was no attempt to control or supervise the use which the church made of this time, and no credits were given for it. Children not so sent remained in the school and were occupied in some manner not counting in their formal studies. The parent then, in signing the church card, chose for his child an elective course in religious subjects, but, in taking this elective the child lost nothing in his formal studies.” This original Gary Plan has undergone some radical changes, most serious among which is the introduction of a cooperative arrangement patterned after the ideals of the Religious Education Association. This affects, of course,
\quellseite{217}
the curriculum and robs all instruction of its confessional character. In the words of the Superintendent of the Community Church Schools, Gary’s goal is “(1) a week-day religious center beside each public school; (2) an opportunity for every child who desires it to have at least one hour of religious instruction during the week; (3) a sufficiently large faculty to provide for small classes; (4) a correlated program of activities and instruction which shall include all the agencies interested in the welfare of the child.” It is evident that not much of positive religious value will be achieved where such a plan is in operation. The growth of the Week Day School movement was quite slow at first, and for several years Gary was the only school in existence; but since 1920 increasing numbers of such schools have been started: 77 in 1920; 131 in 1921, etc.) so that there are over one thousand Week Day schools in operation. It must be said, however, that many schools have passed out of existence soon after their birth and that others diminish in size. \emph{Four types of Week Day schools} are commonly distinguished: (1) the denominational or individual church type which, by the way, is out-stripping all the other types; here the school is governed by the local church and has no connection with any other church or church school in the community; (2) the denominational-cooperating type; here several schools of different denominations cooperate in the form of a loose federation for the promotion of common interests, such as securing public school time, granting credit, advertising, but each church school has its own governing board, its own course of study, and cares for its own children and such others as are not reached by other agencies; (3) the neighborhood or city system; here the schools lose their denominational character completely: the course of study is uniform for all denominations joining in the system, the governing board is made up of representatives of all cooperating churches, and the teachers
\quellseite{218}
are chosen without reference to denominational lines; (4) the pure community type where a board of religious education chosen—not by the churches, but by the community at large, governs the course; at present very few, if any, schools of this type are in operation. It is evident that Lutherans can not take part in any Week Day program similar to types 3 or 4, we can not surrender the choice of subject matter or the selection of teachers, to any agency outside of our own church. Unquestionably Week Day instruction of the individual church type or of the denominational-cooperative type (\# 2) will help remedy the deplorable situation in which Christian education finds itself at the present time. But even though the public school authorities should be found willing to accede to our requests for public school time, one fault would still remain: instruction in the secular branches would continue untouched by the Spirit of Christ; and two views of the world, in conflict with each other, would, during these years of development, compete for mastery in the soul of the child, without the possibility of an adjustment.
